%% =====================================================================
%%  Capítulo 3 -- Fundamentos matemáticos
%%  Formato nuevo (noviembre de 2026). Integra: PrelimMath.tex (2024),
%%  Handout6.tex y Handout6a.tex (versiones anteriores, con la prueba de
%%  existencia y unicidad), y las clases de 2026 de Manuela Chacón
%%  ("Propiedades matemáticas fundamentales": espacio = conjunto +
%%  estructura, ejemplos de métricas, normas de matrices, espacios
%%  completos, jerarquía de regularidad y lema de comparación).
%%  Correcciones respecto a las versiones anteriores: ver \S\ref{sec:pm-cierre}.
%%  Números verificados con herramientas/cap03_numeros.py (43/43) y
%%  codigo/cap03_preliminares/verificarCap03.m (33/33 en Octave 8.4).
%% =====================================================================

\capitulo[cap:preliminares]{Fundamentos matemáticos}
  {Los amores pasados siempre ofenden a los amantes nuevos, por muy muertos que estén aquellos.}
  {Javier Marías}

\begin{objetivos}[partes/03-preliminares.html]
  \item Decir qué estructura tiene un espacio (métrica, norma, producto interno, completitud) y qué preguntas permite responder cada una.
  \item Calcular normas de vectores, de señales y de matrices, y usar la equivalencia de normas en $\R^n$ sin extenderla a espacios de funciones.
  \item Clasificar un campo vectorial como continuo, $\Cont1$, localmente Lipschitz, Lipschitz o globalmente Lipschitz, y calcular una constante de Lipschitz con el jacobiano.
  \item Enunciar y demostrar el teorema de existencia y unicidad local, y decir qué hipótesis falla en cada contraejemplo (no unicidad, escape en tiempo finito).
  \item Acotar soluciones que no se pueden calcular con la desigualdad de Gronwall--Bellman y el lema de comparación.
\end{objetivos}

%% ---------------------------------------------------------------------
\section{¿Por qué matemáticas ahora?}
\label{sec:pm-porque}
%% ---------------------------------------------------------------------

En el capítulo anterior estudiamos $\dot x=f(x)$ con $x\in\R$ suponiendo, casi siempre en silencio, dos cosas: que por cada condición inicial pasa \emph{una} solución, y que esa solución existe \emph{para todo} $t\ge0$. Los dos supuestos fallaron al final del capítulo. La ecuación $\dot x=x^{1/3}$ con $x(0)=0$ tiene infinitas soluciones; $\dot x=x^2$ con $x(0)=1$ tiene una sola, pero deja de existir en $t=1$. Ninguno de los dos es un caso patológico inventado por matemáticos. La no unicidad aparece en el vaciado de un tanque por un orificio (ley de Torricelli, $\dot h=-k\sqrt h$): mirando un tanque vacío no hay manera de saber cuándo se vació, porque hacia atrás en el tiempo la solución no es única. El escape en tiempo finito aparece en cualquier retroalimentación positiva que crezca más rápido que linealmente.

Antes de seguir con sistemas de dimensión dos necesitamos responder, con precisión, tres preguntas sobre el problema de valor inicial
\begin{equation}
  \dot x=f(t,x),\qquad x(t_0)=x_0,\qquad x\in\R^n .
  \label{eq:pm-pvi}
\end{equation}
\begin{enumerate}
  \item ¿Qué condiciones sobre $f$ garantizan que exista una solución y que sea única?
  \item ¿Durante cuánto tiempo existe?
  \item Si no la podemos calcular, ¿podemos al menos \emph{acotarla}?
\end{enumerate}

\notamargen{\textbf{Notación.} El tiempo es el primer argumento: $f(t,x)$, y $f(t,x,u)$ cuando hay entrada. Un sistema es \emph{invariante en el tiempo} si $f$ no depende de $t$ y \emph{variante en el tiempo} si depende. No usamos ``autónomo''.}
Las tres respuestas dependen de poder decir qué tan \emph{cerca} están dos estados, dos soluciones o dos funciones, y qué tan \emph{rápido} puede cambiar $f$ cuando cambia el estado. Lo primero es la geometría de los espacios (métricas, normas, productos internos y completitud, \cref{sec:pm-metricos,sec:pm-normas,sec:pm-producto,sec:pm-completos}). Lo segundo es la regularidad de $f$ (continuidad, diferenciabilidad y la condición de Lipschitz, \cref{sec:pm-regularidad,sec:pm-lipschitz}). Con las dos cosas se prueba el teorema de existencia y unicidad (\cref{sec:pm-existencia}) y se construye la herramienta que más vamos a usar después de este capítulo: el lema de comparación (\cref{sec:pm-comparacion}).

\notamargen{\textbf{Ruta de las clases.} Clase~1: \S\ref{sec:pm-metricos}--\ref{sec:pm-producto}. Clase~2: \S\ref{sec:pm-completos}--\ref{sec:pm-regularidad}. Clase~3: \S\ref{sec:pm-lipschitz} y \S\ref{sec:pm-existencia}. Clase~4: \S\ref{sec:pm-existencia} (pruebas) y \S\ref{sec:pm-comparacion}.}

¿Por qué este capítulo va aquí y no antes? Porque ahora tenemos los ejemplos que lo justifican. Un estudiante que llega a la definición de norma sin haber visto una ecuación con dos soluciones no tiene ninguna razón para quererla; uno que ya vio $\dot x=x^{1/3}$ sabe exactamente qué está buscando.

\begin{ideaclave}
La idea que organiza la primera mitad del capítulo es una sola:
\begin{center}
\sffamily\large\color{snlazul}espacio $=$ conjunto $+$ estructura.
\end{center}
Los elementos pueden ser vectores, funciones del tiempo o incluso palabras. La estructura que se les agrega ---una métrica, una norma, un producto interno--- decide qué preguntas tienen sentido: qué tan lejos, qué tan grande, en qué ángulo, y si las sucesiones que ``deberían'' converger convergen dentro del espacio.
\end{ideaclave}

\begin{figure}[ht]
\centering
\includegraphics[width=\linewidth]{cap03/t_estructuras}
\caption{Cada estructura induce la anterior: un producto interno da una norma y una norma da una métrica. La completitud se agrega por separado; con ella aparecen los espacios de Banach y de Hilbert.}
\label{fig:pm-estructuras}
\end{figure}

\begin{pausa}
Antes de seguir, apueste: si $f$ es continua, ¿el problema \eqref{eq:pm-pvi} tiene una solución única?
\respuesta{No. La continuidad garantiza que existe al menos una solución en un intervalo pequeño (teorema de Peano), pero no que sea única: $\dot x=x^{1/3}$, $x(0)=0$ tiene $x\equiv0$ y $x(t)=(2t/3)^{3/2}$, entre infinitas otras (\cref{ej:pm-nounicidad}). Para la unicidad hace falta la condición de Lipschitz (\cref{thm:pm-local}).}
\end{pausa}

%% ---------------------------------------------------------------------
\section{Espacios métricos: medir distancias}
\label{sec:pm-metricos}
%% ---------------------------------------------------------------------

Tres conjuntos muy distintos van a acompañarnos en esta sección:
\begin{itemize}
  \item $\R^n$, cuyos elementos son vectores $x=[x_1,\dots,x_n]\T$. En un sistema dinámico, el \emph{estado en un instante}.
  \item $C[a,b]$, el conjunto de funciones continuas $f:[a,b]\to\R$. En un sistema dinámico, una \emph{trayectoria completa} o una señal.
  \item Un conjunto de palabras de cuatro letras, $X=\{\text{casa},\text{cama},\text{masa},\text{pasa},\dots\}$.
\end{itemize}
¿Qué tan lejos está un vector de otro? ¿Una función de otra? ¿Una palabra de otra? La respuesta no viene con los objetos: depende de la estructura que se les agregue.

\begin{definicion}[label=def:pm-metrica]{Métrica y espacio métrico}
Sea $X$ un conjunto. Una función $d:X\times X\to\R$ es una \clave{métrica} si para todo $x,y,z\in X$:
\begin{enumerate}[label=(M\arabic*),leftmargin=3em]
  \item $d(x,y)\ge0$ \hfill(no negatividad)
  \item $d(x,y)=0$ si y solo si $x=y$ \hfill(separa puntos)
  \item $d(x,y)=d(y,x)$ \hfill(simetría)
  \item $d(x,y)\le d(x,z)+d(z,y)$ \hfill(desigualdad triangular)
\end{enumerate}
El par $(X,d)$ es un \clave{espacio métrico}.
\end{definicion}

La desigualdad triangular es la que hace el trabajo: dice que no hay atajos, que pasar por un tercer punto nunca acorta el camino. Es la única de las cuatro que hay que probar con cuidado en cada ejemplo.

\begin{ejemplo}[label=ej:pm-metricasRn]{Tres distancias entre los mismos dos puntos}
En $\R^n$ se usan
\[
  d_1(x,y)=\sum_{i=1}^n\abs{x_i-y_i},\qquad
  d_2(x,y)=\Bigl(\sum_{i=1}^n\abs{x_i-y_i}^2\Bigr)^{1/2},\qquad
  d_\infty(x,y)=\max_i\abs{x_i-y_i},
\]
y, en general, $d_p(x,y)=\bigl(\sum_i\abs{x_i-y_i}^p\bigr)^{1/p}$ con $1\le p<\infty$ ($d_\infty$ es el límite de $d_p$ cuando $p\to\infty$). Calcule las tres para $x=[1,1]\T$ y $y=[4,3]\T$, y pruebe que $d_1$ es una métrica.
\solucion $d_1=3+2=5$ (el taxista que solo se mueve por calles), $d_2=\sqrt{9+4}=\sqrt{13}\approx3.61$ (la línea recta) y $d_\infty=\max\{3,2\}=3$ (la coordenada que más difiere). Los puntos son los mismos; la distancia cambia con la métrica (\cref{fig:pm-distancias}, izquierda). Para $d_1$, (M1)--(M3) son inmediatas. Para (M4), en cada componente $x_i-y_i=(x_i-z_i)+(z_i-y_i)$ y la desigualdad triangular de los reales da $\abs{x_i-y_i}\le\abs{x_i-z_i}+\abs{z_i-y_i}$; sumando sobre $i$ se obtiene $d_1(x,y)\le d_1(x,z)+d_1(z,y)$.
\end{ejemplo}

\begin{ejemplo}[label=ej:pm-hamming]{Una distancia entre palabras}
En el conjunto de palabras de cuatro letras, sea $d(u,v)$ el número de posiciones en las que $u$ y $v$ difieren (distancia de Hamming). Calcule $d(\text{casa},\text{cama})$, $d(\text{casa},\text{masa})$ y $d(\text{cama},\text{masa})$, y diga por qué es una métrica.
\solucion $d(\text{casa},\text{cama})=1$, $d(\text{casa},\text{masa})=1$ y $d(\text{cama},\text{masa})=2$. Es una métrica: (M1)--(M3) son claras, y para (M4), si $u$ y $v$ difieren en la posición $i$, entonces $w$ difiere de $u$ o de $v$ en esa posición; cada posición que cuenta en $d(u,v)$ cuenta al menos una vez en $d(u,w)+d(w,v)$. La misma distancia mide cuántos bits cambió el ruido en un mensaje digital: los objetos no tienen que ser números para que tenga sentido hablar de cercanía.
\end{ejemplo}

\script{cap03_preliminares/C3_Ej1_Metricas.m}
\begin{figure}[ht]
\centering
\includegraphics[width=\linewidth]{cap03/f_distancias}
\caption{Izquierda: tres métricas en $\R^2$ para los mismos dos puntos. Derecha: entre funciones de $C[0,1]$, $d_1$ es el área entre las curvas y $d_\infty$ la mayor separación vertical.}
\label{fig:pm-distancias}
\end{figure}

Entre funciones continuas $f,g\in C[a,b]$ se definen, de la misma manera,
\begin{equation}
  d_p(f,g)=\Bigl(\int_a^b\abs{f(\tau)-g(\tau)}^p\,\dif\tau\Bigr)^{1/p},\quad 1\le p<\infty,\qquad
  d_\infty(f,g)=\max_{a\le t\le b}\abs{f(t)-g(t)} .
  \label{eq:pm-dfunciones}
\end{equation}
La suma sobre las componentes se reemplaza por una integral sobre el tiempo. Para $f(t)=\sin2\pi t$ y $g(t)=t$ en $[0,1]$, $d_1(f,g)=0.921$ y $d_\infty(f,g)=1.763$, alcanzada en $t=0.775$ (\cref{fig:pm-distancias}, derecha). Que $d_p(f,g)=0$ implique $f=g$ usa la continuidad: una función continua y no negativa con integral nula es idénticamente cero.

\codigorepo{cap03_preliminares/C3_Ej1_Metricas.m}{vectores}{Tres métricas para los mismos dos puntos}

\begin{trampa}[label=tr:pm-parece]{Parecer una distancia no basta}
\plausible{$d(x,y)=(x-y)^2$ es no negativa, simétrica y vale cero solo si $x=y$; por lo tanto es una métrica en $\R$.}
Falta la desigualdad triangular, y falla: $d(0,2)=4$, pero $d(0,1)+d(1,2)=1+1=2$. Pasar por $z=1$ ``acorta'' el camino. Elevar al cuadrado castiga las distancias grandes más que proporcionalmente, y eso es incompatible con (M4).
\correcto{Para concluir que una función es una métrica hay que verificar las cuatro propiedades; la triangular es la que casi siempre decide. En cambio, $\abs{x-y}^{1/2}$ sí es una métrica en $\R$ (pruébelo: $\sqrt{a+b}\le\sqrt a+\sqrt b$).}
\end{trampa}

\codigorepo{cap03_preliminares/C3_Ej1_Metricas.m}{nometrica}{$(x-y)^2$ no es métrica}

\begin{pausa}
En un tablero de ajedrez, ¿qué métrica cuenta el número de movidas que necesita un rey para ir de una casilla a otra? ¿Y una torre que solo puede moverse una casilla a la vez?
\respuesta{El rey se mueve una casilla en cualquier dirección, incluidas las diagonales: necesita $\max\{\abs{\Delta x},\abs{\Delta y}\}$ movidas, la métrica $d_\infty$. La torre restringida a una casilla por movida necesita $\abs{\Delta x}+\abs{\Delta y}$, la métrica $d_1$.}
\end{pausa}

%% ---------------------------------------------------------------------
\section{Espacios normados: medir tamaños}
\label{sec:pm-normas}
%% ---------------------------------------------------------------------

Una métrica responde ``¿qué tan lejos están $x$ y $y$?''. En un espacio vectorial ---donde se puede sumar y multiplicar por escalares--- tiene sentido una pregunta más básica: ``¿qué tan grande es $x$?''.

\begin{definicion}[label=def:pm-norma]{Norma y espacio normado}
Sea $X$ un espacio vectorial real. Una función $\norm{\cdot}:X\to\R$ es una \clave{norma} si para todo $x,y\in X$ y todo $\alpha\in\R$:
\begin{enumerate}[label=(N\arabic*),leftmargin=3em]
  \item $\norm x\ge0$, y $\norm x=0$ si y solo si $x=0$ \hfill(definida positiva)
  \item $\norm{\alpha x}=\abs\alpha\,\norm x$ \hfill(homogeneidad)
  \item $\norm{x+y}\le\norm x+\norm y$ \hfill(desigualdad triangular)
\end{enumerate}
El par $(X,\norm{\cdot})$ es un \clave{espacio normado}.
\end{definicion}

Toda norma induce una métrica, $d(x,y)\defeq\norm{x-y}$: la distancia entre dos elementos es el tamaño del desplazamiento que lleva de uno al otro. Las propiedades (M1)--(M4) salen de (N1)--(N3); por ejemplo, $\norm{x-y}=\norm{(x-z)+(z-y)}\le\norm{x-z}+\norm{z-y}$. Lo contrario no es cierto: la métrica de Hamming no viene de ninguna norma (las palabras no se suman).

En $\R^n$, las métricas $d_p$ de la sección anterior vienen de las \clave{normas-$p$}
\begin{equation}
  \norm x_p=\Bigl(\sum_{i=1}^n\abs{x_i}^p\Bigr)^{1/p}\quad(1\le p<\infty),\qquad \norm x_\infty=\max_i\abs{x_i} .
  \label{eq:pm-normasp}
\end{equation}
Para $x=[3,4]\T$: $\norm x_1=7$, $\norm x_2=5$ y $\norm x_\infty=4$. Cuando no se indica el subíndice, $\norm x$ es la norma euclidiana $\norm x_2$.

\script{cap03_preliminares/C3_Ej2_Normas.m}
\begin{figure}[ht]
\centering
\includegraphics[width=\linewidth]{cap03/f_bolas}
\caption{Izquierda: la ``esfera unitaria'' $\{x:\norm x_p=1\}$ es un rombo para $p=1$, un círculo para $p=2$ y un cuadrado para $p=\infty$. Para $p=1/2$ la figura no es convexa: no es una norma. Derecha: en $C[0,1]$ la sucesión $t^k$ tiene norma del máximo $1$ y norma $\mathcal{L}_1$ que tiende a cero.}
\label{fig:pm-bolas}
\end{figure}

\begin{trampa}[label=tr:pm-pmenor]{La fórmula no hace la norma}
\plausible{$\norm x_p=\bigl(\sum_i\abs{x_i}^p\bigr)^{1/p}$ es una norma para todo $p>0$; para $p=1/2$ simplemente pesa más las componentes pequeñas.}
Para $p<1$ falla la desigualdad triangular. Con $a=[1,0]\T$, $b=[0,1]\T$ y $p=1/2$: $\norm a=\norm b=1$, pero $\norm{a+b}=(1+1)^2=4>2$. Geométricamente, la ``bola unitaria'' de la \cref{fig:pm-bolas} no es convexa: el segmento entre dos puntos de la bola se sale de ella.
\correcto{$\norm\cdot_p$ es una norma si y solo si $p\ge1$ (la triangular para $p\ge1$ es la desigualdad de Minkowski). La bola unitaria de cualquier norma es convexa y simétrica respecto al origen.}
\end{trampa}

\codigorepo{cap03_preliminares/C3_Ej2_Normas.m}{pmedio}{$p=1/2$ viola la desigualdad triangular}

%% ---------------------------------------------------------------------
\subsection{Equivalencia de normas en \texorpdfstring{$\R^n$}{Rn}}
%% ---------------------------------------------------------------------

Los valores cambian con la norma, pero en $\R^n$ nunca se alejan demasiado entre sí.

\begin{teorema}[label=thm:pm-equivalencia]{Equivalencia de normas en dimensión finita}
Si $\norm\cdot_\alpha$ y $\norm\cdot_\beta$ son dos normas en $\R^n$, existen constantes $c_1,c_2>0$ tales que
\[
  c_1\norm x_\alpha\le\norm x_\beta\le c_2\norm x_\alpha\qquad\text{para todo }x\in\R^n .
\]
En particular,
\begin{equation}
  \norm x_2\le\norm x_1\le\sqrt n\,\norm x_2,\qquad
  \norm x_\infty\le\norm x_2\le\sqrt n\,\norm x_\infty,\qquad
  \norm x_\infty\le\norm x_1\le n\,\norm x_\infty .
  \label{eq:pm-equiv}
\end{equation}
\end{teorema}

\begin{proof}
Las tres desigualdades de \eqref{eq:pm-equiv}. Sea $m=\max_i\abs{x_i}=\norm x_\infty$. (i) $m\le\sum_i\abs{x_i}\le n\,m$. (ii) $m^2\le\sum_ix_i^2\le n\,m^2$; tome raíz. (iii) $\norm x_1^2=\sum_i\abs{x_i}^2+\sum_{i\ne j}\abs{x_i}\abs{x_j}\ge\norm x_2^2$; y, por Cauchy--Schwarz (\cref{thm:pm-cs}) con el vector $[1,\dots,1]\T$, $\norm x_1=\sum_i1\cdot\abs{x_i}\le\sqrt n\,\norm x_2$. Las cotas son ajustadas: $x=[1,0,\dots,0]\T$ da la igualdad a la izquierda y $x=[1,1,\dots,1]\T$ a la derecha.
El caso general: basta compararlas con $\norm\cdot_2$. Toda norma es continua respecto a $\norm\cdot_2$, porque $\abs{\norm x_\alpha-\norm y_\alpha}\le\norm{x-y}_\alpha\le\bigl(\max_i\norm{e_i}_\alpha\bigr)\norm{x-y}_1\le\bigl(\max_i\norm{e_i}_\alpha\bigr)\sqrt n\norm{x-y}_2$. La esfera $\{\norm x_2=1\}$ es compacta (\cref{sec:pm-conjuntos}), así que $\norm\cdot_\alpha$ alcanza en ella un mínimo $c_1>0$ (no es cero porque $x\neq0$) y un máximo $c_2$. Por homogeneidad, $c_1\norm x_2\le\norm x_\alpha\le c_2\norm x_2$ para todo $x$.
\end{proof}

¿Para qué sirve? Para elegir la norma que haga más fácil una cuenta. Si se prueba que $\norm{x(t)}_\infty\to0$, también $\norm{x(t)}_2\to0$; si se prueba que $\norm{x(t)}_1\le ce^{-t}$, entonces $\norm{x(t)}_2\le ce^{-t}$. Las propiedades \emph{cualitativas} ---acotado, convergente, estable--- no dependen de la norma en $\R^n$; las constantes sí. Con $x=[3,4]\T$: $5\le7\le5\sqrt2\approx7.07$.

\begin{cuidado}[Supremo y máximo]
El \emph{supremo} de un conjunto $S\subset\R$ es su menor cota superior; el \emph{máximo} es un elemento de $S$ que es cota superior. $S_1=[0,1]$ tiene $\sup S_1=\max S_1=1$; $S_2=(0,1)$ tiene $\sup S_2=1$ y no tiene máximo. Por eso las normas de señales se escriben con $\sup$: la señal $1-e^{-t}$, $t\ge0$, nunca llega a $1$, y su norma del máximo es $1$.
\end{cuidado}

%% ---------------------------------------------------------------------
\subsection{Normas de señales}
%% ---------------------------------------------------------------------

Una trayectoria $x:[0,\infty)\to\R^n$ es un elemento de un espacio de funciones. Se puede medir el tamaño del estado en un instante, $\norm{x(t)}$, o el tamaño de \emph{toda} la trayectoria:
\begin{equation}
  \norm x_{\mathcal{L}_p}=\Bigl(\int_0^\infty\norm{x(t)}^p\,\dif t\Bigr)^{1/p}\quad(1\le p<\infty),\qquad
  \norm x_{\Linf}=\sup_{t\ge0}\norm{x(t)} .
  \label{eq:pm-Lp}
\end{equation}
\notamargen{Para funciones medibles generales, $\mathcal{L}_p$ se define con la integral de Lebesgue y sus elementos son clases de funciones que coinciden casi en todas partes. En el curso basta pensar en señales continuas a trozos.}
Se dice que $x\in\mathcal{L}_p$ si esa cantidad es \emph{finita}. $\norm x_{\mathcal{L}_1}$ es la magnitud acumulada, $\norm x_{\Ldos}^2$ es una \emph{energía} (si $x$ es una corriente por un resistor de $1\,\Omega$, es la energía disipada) y $\norm x_{\Linf}$ es el pico. Pertenecer a uno de estos espacios no implica pertenecer a otro (\cref{fig:pm-senales}): $1/(1+t)\in\Ldos\cap\Linf$ pero $\notin\mathcal{L}_1$; $\sin t\in\Linf$ solamente; $t^{-1/2}$ en $(0,1)$ (y cero después) está en $\mathcal{L}_1$ pero no en $\Ldos$ ni en $\Linf$. Estas normas vuelven en el capítulo~11, donde la energía de una señal es la que entra por un puerto.

\begin{figure}[ht]
\centering
\includegraphics[width=\linewidth]{cap03/f_senales}
\caption{Cuatro señales y los espacios $\mathcal{L}_p[0,\infty)$ a los que pertenecen. Una señal acotada no tiene por qué tener energía finita, y una con área finita no tiene por qué ser acotada.}
\label{fig:pm-senales}
\end{figure}

\begin{trampa}[label=tr:pm-dimension]{La equivalencia de normas es un privilegio de $\R^n$}
\plausible{Todas las normas son equivalentes; si una señal es pequeña en $\mathcal{L}_1$, es pequeña en la norma del máximo.}
En espacios de funciones es falso. En $C[0,1]$, $f_k(t)=t^k$ cumple $\norm{f_k}_\infty=1$ para todo $k$, mientras que $\norm{f_k}_{\mathcal{L}_1}=1/(k+1)\to0$ (\cref{fig:pm-bolas}, derecha). No existe $c$ con $\norm f_\infty\le c\norm f_{\mathcal{L}_1}$: la sucesión ``converge a cero'' en una norma y no en la otra. La prueba del \cref{thm:pm-equivalencia} usó que la esfera unitaria de $\R^n$ es compacta, y eso falla en dimensión infinita.
\correcto{En $\R^n$ se puede cambiar de norma libremente (con constantes). Con señales hay que decir \emph{en qué norma} es pequeña una señal o converge una sucesión, porque las conclusiones cambian.}
\end{trampa}

\codigorepo{cap03_preliminares/C3_Ej2_Normas.m}{funciones}{$\norm{t^k}_\infty=1$ mientras $\norm{t^k}_{\mathcal{L}_1}\to0$}

%% ---------------------------------------------------------------------
\subsection{Normas de matrices}
%% ---------------------------------------------------------------------

Una matriz $A\in\R^{m\times n}$ transforma vectores: $x\mapsto Ax$. Su tamaño natural es cuánto puede amplificar.

\begin{definicion}[label=def:pm-inducida]{Norma inducida de una matriz}
\[
  \norm A_p=\sup_{x\neq0}\frac{\norm{Ax}_p}{\norm x_p}=\max_{\norm x_p=1}\norm{Ax}_p .
\]
Para $p=1,2,\infty$ hay fórmulas cerradas:
\[
  \norm A_1=\max_j\sum_{i=1}^m\abs{a_{ij}}\ \ \text{(columnas)},\qquad
  \norm A_2=\sqrt{\lmax(A\T A)},\qquad
  \norm A_\infty=\max_i\sum_{j=1}^n\abs{a_{ij}}\ \ \text{(filas)} .
\]
\end{definicion}

\begin{proposicion}[label=prop:pm-matrices]{Propiedades de las normas inducidas}
Para $A\in\R^{m\times n}$ y $B\in\R^{n\times\ell}$:
\begin{enumerate}
  \item $\norm{Ax}_p\le\norm A_p\norm x_p$ \quad(es la definición: ninguna dirección amplifica más);
  \item $\norm{AB}_p\le\norm A_p\norm B_p$ \quad(submultiplicativa);
  \item $\tfrac1{\sqrt n}\norm A_\infty\le\norm A_2\le\sqrt m\,\norm A_\infty$ y $\tfrac1{\sqrt m}\norm A_1\le\norm A_2\le\sqrt n\,\norm A_1$;
  \item $\norm A_2\le\sqrt{\norm A_1\norm A_\infty}$.
\end{enumerate}
\end{proposicion}

\begin{ejemplo}[label=ej:pm-matriz]{Tres normas de la misma matriz}
Calcule $\norm A_1$, $\norm A_\infty$ y $\norm A_2$ para $A=\begin{bmatrix}1&2\\-3&4\end{bmatrix}$, y compare $\norm A_2$ con el mayor $\abs{\lambda}$.
\solucion Columnas: $1+3=4$ y $2+4=6$, luego $\norm A_1=6$. Filas: $1+2=3$ y $3+4=7$, luego $\norm A_\infty=7$. $A\T A=\begin{bmatrix}10&-10\\-10&20\end{bmatrix}$ tiene valores propios $15\pm5\sqrt5$, así que $\norm A_2=\sqrt{15+5\sqrt5}\approx5.117\le\sqrt{6\cdot7}\approx6.48$. Los valores propios de $A$ son $\frac{5\pm j\sqrt{15}}{2}$, con módulo $\sqrt{10}\approx3.16$: el \emph{radio espectral} $\rho(A)=\max\abs{\lambda_i}$ es menor que $\norm A_2$. La \cref{fig:pm-matriz} muestra por qué: el círculo unitario se transforma en una elipse cuyo semieje mayor es $\norm A_2$.
\end{ejemplo}

\codigorepo{cap03_preliminares/C3_Ej3_NormasMatrices.m}{inducidas}{Las tres normas inducidas, con fórmula y con \cmd{norm}}

\begin{figure}[ht]
\centering
\includegraphics[width=\linewidth]{cap03/f_matriz}
\caption{Izquierda: $A$ convierte el círculo unitario en una elipse; el semieje mayor es $\norm A_2$, mayor que $\rho(A)$ (círculo discontinuo). Derecha: la matriz nilpotente $N$ tiene $\rho(N)=0$, pero aplasta el círculo en un segmento de longitud $2$: $\norm N_2=1$.}
\label{fig:pm-matriz}
\end{figure}

\begin{trampa}[label=tr:pm-radio]{La norma 2 no es el mayor valor propio}
\plausible{$\norm A_2=\max_i\abs{\lambda_i(A)}$: la amplificación máxima la da el valor propio más grande.}
Solo es cierto si $A$ es normal ($A\T A=AA\T$), por ejemplo simétrica. En general $\rho(A)\le\norm A$ para \emph{cualquier} norma inducida (si $Av=\lambda v$, entonces $\abs\lambda\norm v=\norm{Av}\le\norm A\norm v$), pero la desigualdad puede ser estricta y enorme. $N=\begin{bmatrix}0&1\\0&0\end{bmatrix}$ tiene $\rho(N)=0$ y $\norm N_2=1$.
\correcto{Use $\norm A_2=\sqrt{\lmax(A\T A)}$, el mayor valor singular. La diferencia importa en estabilidad: una matriz Hurwitz puede amplificar mucho antes de decaer ($\norm{e^{At}}$ crece transitoriamente si $A$ no es normal), y las funciones de Lyapunov cuadráticas del capítulo~7 son la manera correcta de manejarlo.}
\end{trampa}

\codigorepo{cap03_preliminares/C3_Ej3_NormasMatrices.m}{espectral}{Radio espectral y norma 2}

%% ---------------------------------------------------------------------
\section{Producto interno: medir ángulos}
\label{sec:pm-producto}
%% ---------------------------------------------------------------------

Una norma dice qué tan grande es un vector, pero no si dos vectores ``apuntan para el mismo lado''. Esa información ---ángulos, ortogonalidad, proyecciones--- la da el producto interno.

\begin{definicion}[label=def:pm-pinterno]{Producto interno}
Sea $X$ un espacio vectorial real. Una función $\inner{\cdot}{\cdot}:X\times X\to\R$ es un \clave{producto interno} si para todo $x,y,z\in X$ y $\alpha\in\R$:
\begin{enumerate}[label=(P\arabic*),leftmargin=3em]
  \item $\inner{x+y}{z}=\inner xz+\inner yz$ y $\inner{\alpha x}{y}=\alpha\inner xy$ \hfill(lineal en el primer argumento)
  \item $\inner xy=\inner yx$ \hfill(simétrico)
  \item $\inner xx>0$ si $x\neq0$ \hfill(definido positivo)
\end{enumerate}
\end{definicion}

En $\R^n$ el producto interno usual es $\inner xy=x\T y=\sum_ix_iy_i$, y en $C[a,b]$ es $\inner fg=\int_a^bf(t)g(t)\,\dif t$. Un producto interno induce una norma, $\norm x=\sqrt{\inner xx}$: el de $\R^n$ induce $\norm\cdot_2$ y el de $C[a,b]$ induce $\norm\cdot_{\Ldos}$. Hay, por lo tanto, una cadena
\[
  \text{producto interno}\implica\text{norma}\implica\text{métrica}.
\]
Que $\sqrt{\inner xx}$ cumpla la desigualdad triangular no es obvio; es consecuencia del resultado central de esta sección.

\begin{teorema}[label=thm:pm-cs]{Desigualdad de Cauchy--Schwarz}
En un espacio con producto interno, con $\norm x=\sqrt{\inner xx}$,
\[
  \abs{\inner xy}\le\norm x\,\norm y\qquad\text{para todo }x,y,
\]
con igualdad si y solo si $x$ y $y$ son linealmente dependientes.
\end{teorema}

\begin{proof}
Si $y=0$ no hay nada que probar. Si $y\neq0$, para todo $s\in\R$,
\[
  0\le\inner{x-sy}{x-sy}=\norm x^2-2s\inner xy+s^2\norm y^2 .
\]
Es una parábola en $s$ que nunca es negativa, así que su discriminante no es positivo: $4\inner xy^2-4\norm x^2\norm y^2\le0$. La igualdad ocurre si y solo si la parábola toca el cero, es decir, si $x=sy$ para algún $s$.
\end{proof}

Con Cauchy--Schwarz, la triangular de la norma inducida es una línea: $\norm{x+y}^2=\norm x^2+2\inner xy+\norm y^2\le(\norm x+\norm y)^2$. Y tiene sentido definir el \clave{ángulo} entre $x$ y $y$ como el $\theta\in[0,\pi]$ con $\cos\theta=\inner xy/(\norm x\norm y)$: el cociente está en $[-1,1]$ \emph{gracias} al teorema. El signo de $\inner xy$ dice si el ángulo es agudo ($>0$), recto ($=0$, vectores \clave{ortogonales}) u obtuso ($<0$). Esa lectura geométrica es la que se usa en Lyapunov: $\Vdot=\nabla V(x)\T f(x)<0$ dice que el campo forma un ángulo obtuso con el gradiente de $V$, es decir, que apunta ``cuesta abajo''.

\begin{trampa}[label=tr:pm-circular]{Cauchy--Schwarz no sale de $\abs{\cos\theta}\le1$}
\plausible{$\inner xy=\norm x\norm y\cos\theta$ y $\abs{\cos\theta}\le1$; por lo tanto $\abs{\inner xy}\le\norm x\norm y$.}
En $\R^2$ y $\R^3$ el argumento es correcto porque el ángulo se define con la geometría euclidiana. Pero en un espacio general ---funciones, polinomios, $\R^{100}$--- el ángulo \emph{se define} con la fórmula $\cos\theta=\inner xy/(\norm x\norm y)$, y para que esa definición tenga sentido hace falta saber antes que el cociente está en $[-1,1]$. El argumento es circular.
\correcto{La prueba general usa solo los axiomas (la parábola del \cref{thm:pm-cs}). La figura en $\R^2$ es una buena intuición, no una demostración.}
\end{trampa}

\begin{teorema}[label=thm:pm-holder]{Desigualdad de Hölder}
Sean $p,q\in[1,\infty]$ con $\frac1p+\frac1q=1$ (exponentes conjugados; $p=1$ va con $q=\infty$). Para $x,y\in\R^n$ y para funciones $f\in\mathcal{L}_p$, $g\in\mathcal{L}_q$,
\[
  \sum_{i=1}^n\abs{x_iy_i}\le\norm x_p\norm y_q,\qquad
  \int\abs{f(t)g(t)}\,\dif t\le\norm f_{\mathcal{L}_p}\norm g_{\mathcal{L}_q} .
\]
\end{teorema}

Con $p=q=2$ se recupera Cauchy--Schwarz. Hölder permite acotar un producto usando normas \emph{distintas} para cada factor; por ejemplo, $\abs{x\T y}\le\norm x_1\norm y_\infty$, que se usa cuando una señal es acotada y la otra tiene área finita: la potencia que entra a un sistema con una entrada acotada y una salida en $\mathcal{L}_1$ tiene energía total finita.

\begin{pausa}
¿Proviene $\norm\cdot_1$ de algún producto interno en $\R^2$? Pista: si $\norm x=\sqrt{\inner xx}$, entonces $\norm{x+y}^2+\norm{x-y}^2=2\norm x^2+2\norm y^2$ (ley del paralelogramo).
\respuesta{No. Con $x=[1,0]\T$, $y=[0,1]\T$: $\norm{x+y}_1^2+\norm{x-y}_1^2=4+4=8$, mientras que $2\norm x_1^2+2\norm y_1^2=4$. La ley del paralelogramo falla, así que ningún producto interno induce $\norm\cdot_1$. Entre las normas-$p$, solo $p=2$ viene de un producto interno: por eso la geometría de ángulos y proyecciones se hace con la norma euclidiana.}
\end{pausa}

%% ---------------------------------------------------------------------
\section{Espacios completos: Banach y Hilbert}
\label{sec:pm-completos}
%% ---------------------------------------------------------------------

La prueba de existencia de soluciones (\cref{sec:pm-existencia}) construye una sucesión de funciones que se acercan unas a otras y afirma que tiene un límite. Que ese límite exista \emph{dentro del espacio} no es automático. Es la propiedad que falta.

\begin{definicion}[label=def:pm-cauchy]{Sucesión de Cauchy y espacio completo}
En un espacio métrico $(X,d)$, una sucesión $\{x_k\}$ es \clave{de Cauchy} si para todo $\epsilon>0$ existe $N$ tal que $d(x_k,x_m)<\epsilon$ para todo $k,m\ge N$. El espacio es \clave{completo} si toda sucesión de Cauchy converge a un elemento \emph{de} $X$.
\end{definicion}

Toda sucesión convergente es de Cauchy (si $x_k\to x$, los términos se acercan a $x$ y, por la triangular, entre sí). Lo contrario es lo que define la completitud. Los racionales no son completos: $1,\ 1.4,\ 1.41,\ 1.414,\dots$ es de Cauchy en $\mathbb{Q}$ y su límite, $\sqrt2$, no es racional. Un espacio completo es uno sin ``huecos''.

\begin{definicion}[label=def:pm-banach]{Espacios de Banach y de Hilbert}
Un espacio normado que es completo con la métrica $d(x,y)=\norm{x-y}$ es un \clave{espacio de Banach}. Un espacio con producto interno que es completo con la norma inducida es un \clave{espacio de Hilbert}.
\end{definicion}

Todo espacio de Hilbert es de Banach (su norma inducida lo hace normado y completo); no todo Banach es Hilbert: $(\R^n,\norm\cdot_1)$ es Banach y su norma no viene de un producto interno.

\begin{ejemplo}[label=ej:pm-Rnbanach]{$\R^n$ es completo con cualquier norma-$p$}
Pruebe que $(\R^n,\norm\cdot_p)$, $1\le p\le\infty$, es un espacio de Banach.
\solucion Sea $\{x^{(k)}\}$ de Cauchy. Como $\abs{x^{(k)}_i-x^{(m)}_i}\le\norm{x^{(k)}-x^{(m)}}_\infty\le\norm{x^{(k)}-x^{(m)}}_p$, cada componente $\{x_i^{(k)}\}$ es de Cauchy en $\R$. Los reales son completos, así que $x_i^{(k)}\to x_i^*$ para cada $i$. Sea $x^*=[x_1^*,\dots,x_n^*]\T$. Convergencia componente a componente da $\norm{x^{(k)}-x^*}_\infty\to0$ y, por la equivalencia de normas, $\norm{x^{(k)}-x^*}_p\to0$. Con el producto interno usual, $\R^n$ es además un espacio de Hilbert.
\end{ejemplo}

El espacio de funciones que vamos a usar para probar la existencia de soluciones es $C[a,b]$ (funciones continuas $x:[a,b]\to\R^n$) con la \emph{norma del máximo}
\begin{equation}
  \norm x_C=\max_{a\le t\le b}\norm{x(t)} .
  \label{eq:pm-normaC}
\end{equation}
$(C[a,b],\norm\cdot_C)$ es un espacio de Banach: si $\{x_k\}$ es de Cauchy en esa norma, para cada $t$ la sucesión $\{x_k(t)\}$ es de Cauchy en $\R^n$ y tiene límite $x(t)$; la convergencia es \emph{uniforme} en $t$, y un límite uniforme de funciones continuas es continuo. El límite vive en $C[a,b]$.

\begin{trampa}[label=tr:pm-completo]{Producto interno no es lo mismo que Hilbert}
\plausible{$C[0,1]$ con $\inner fg=\int_0^1fg\,\dif t$ tiene producto interno; por lo tanto es un espacio de Hilbert.}
Falta la completitud, y falla. Las rampas continuas $f_k(t)$, que valen $0$ hasta $t=\tfrac12$, suben con pendiente $k$ y valen $1$ después, forman una sucesión de Cauchy en $\Ldos$: $\norm{f_k-f_m}_{\Ldos}\to0$ (por ejemplo, $\norm{f_{10}-f_{20}}_{\Ldos}=1/\sqrt{120}\approx0.091$). Pero se acercan al escalón, que no es continuo (\cref{fig:pm-completitud}). En la norma del máximo la misma sucesión ni siquiera es de Cauchy: $\norm{f_k-f_{2k}}_C=\tfrac12$ para todo $k$.
\correcto{$(C[0,1],\norm\cdot_C)$ es Banach; $(C[0,1],\norm\cdot_{\Ldos})$ no es completo. Su completación ---agregar todos los límites--- es el espacio $\Ldos[0,1]$, que sí es de Hilbert. Por eso la prueba de existencia trabaja con la norma del máximo.}
\end{trampa}

\script{cap03_preliminares/C3_Ej4_Completitud.m}
\begin{figure}[ht]
\centering
\includegraphics[width=0.7\linewidth]{cap03/f_completitud}
\caption{Rampas continuas que, en la norma $\Ldos$, se acercan entre sí y al escalón: $\norm{f_k-\text{escalón}}_{\Ldos}=1/\sqrt{3k}$. El límite no está en $C[0,1]$.}
\label{fig:pm-completitud}
\end{figure}

\codigorepo{cap03_preliminares/C3_Ej4_Completitud.m}{rampas}{Cauchy en $\Ldos$, no en la norma del máximo}

\begin{historia}[Banach, Hilbert y el espacio $\Ldos$]
Hilbert estudió entre 1904 y 1910 ecuaciones integrales en lo que hoy llamamos el espacio $\ell^2$ de sucesiones de cuadrado sumable; la estructura abstracta de espacio de Hilbert la axiomatizó von Neumann hacia 1929. Stefan Banach definió los espacios normados completos en su tesis, publicada en 1922, donde también aparece el principio de contracción que usaremos en la \cref{sec:pm-existencia} \citep{Ban22}. Que $\Ldos$ es completo es el teorema de Riesz--Fischer (1907).
\end{historia}

%% ---------------------------------------------------------------------
\section{Conjuntos en \texorpdfstring{$\R^n$}{Rn}: el vocabulario mínimo}
\label{sec:pm-conjuntos}
%% ---------------------------------------------------------------------

Los teoremas de este capítulo y de los siguientes piden que un conjunto sea abierto, cerrado, compacto, convexo o un dominio. Estas son las definiciones, en $\R^n$ con cualquier norma (por la equivalencia de normas, no importa cuál).

\begin{definicion}[label=def:pm-conjuntos]{Abierto, cerrado, acotado, compacto}
Sea $S\subseteq\R^n$ y sea $\{x:\norm{x-p}<\delta\}$ la \clave{vecindad} de radio $\delta$ de $p$.
\begin{itemize}
  \item $S$ es \clave{abierto} si cada $p\in S$ tiene una vecindad contenida en $S$.
  \item $p$ es un \clave{punto límite} de $S$ si toda vecindad de $p$ contiene un punto de $S$ distinto de $p$. $S$ es \clave{cerrado} si contiene todos sus puntos límite; equivalentemente, si el límite de toda sucesión convergente de puntos de $S$ está en $S$.
  \item $S$ es \clave{acotado} si existe $r$ con $\norm x\le r$ para todo $x\in S$.
  \item $S$ es \clave{compacto} si es cerrado y acotado.
\end{itemize}
\end{definicion}

\notamargen{En espacios métricos generales, ``compacto'' significa que toda sucesión tiene una subsucesión convergente en $S$. En $\R^n$ eso equivale a cerrado y acotado (Heine--Borel). En dimensión infinita no: la bola unitaria cerrada de $C[0,1]$ es cerrada y acotada y no es compacta.}
La bola cerrada $\Bola{r}=\{x:\norm x\le r\}$ es compacta; la abierta $\{\norm x<r\}$ es abierta y acotada, pero no cerrada. La compacidad es lo que garantiza que una función continua alcance su máximo y su mínimo, y que toda sucesión tenga una subsucesión convergente; las dos cosas se usan en las pruebas de existencia global (\cref{thm:pm-compacto}) y de Poincaré--Bendixson (capítulo~5).

\begin{definicion}[label=def:pm-convexo]{Convexo, conexo, dominio}
$W\subseteq\R^n$ es \clave{convexo} si contiene el segmento entre dos cualesquiera de sus puntos: $\lambda x+(1-\lambda)y\in W$ para todo $x,y\in W$ y $\lambda\in[0,1]$. Un abierto es \clave{conexo} si dos puntos cualesquiera se pueden unir con una poligonal contenida en él. Un \clave{dominio} es un conjunto abierto y conexo.
\end{definicion}

\begin{figure}[ht]
\centering
\includegraphics[width=\linewidth]{cap03/f_conjuntos}
\caption{De izquierda a derecha: un abierto (el borde discontinuo no le pertenece), un compacto, un conjunto no convexo y un abierto que no es conexo.}
\label{fig:pm-conjuntos}
\end{figure}

La convexidad aparece en un solo lugar, pero es esencial: para acotar $f(x)-f(y)$ con el jacobiano hay que recorrer el segmento entre $x$ y $y$, y ese segmento tiene que quedarse donde la cota vale (\cref{lem:pm-jacobiano}). Los dominios aparecen en los enunciados porque en un abierto se puede derivar en todas las direcciones, y porque en un conjunto conexo una propiedad local se puede ``propagar''.

%% ---------------------------------------------------------------------
\section{Regularidad del campo vectorial}
\label{sec:pm-regularidad}
%% ---------------------------------------------------------------------

Casi nunca se analiza la ecuación de la planta tal como viene, $\dot x=f_p(t,x,u)$, con $x(t)\in\R^n$ y $u(t)\in\R^m$. Se le conecta una ley de control $u=\gamma(t,x)$ y se analiza el lazo cerrado
\begin{equation}
  \dot x=f_p\bigl(t,x,\gamma(t,x)\bigr)\eqqcolon f(t,x) .
  \label{eq:pm-lazo}
\end{equation}
\notamargen{Una entrada $u(t)$ conocida (una referencia, una perturbación medida) también entra a través de la dependencia en $t$. Si $u$ tiene saltos, $f$ tiene saltos en $t$: por eso se pide continuidad a trozos.}
Las propiedades de $f$ heredan las de $f_p$ y las de $\gamma$: un controlador discontinuo (un relé, una saturación con conmutación) produce un $f$ discontinuo. Esta sección fija el vocabulario para decir qué tan ``suave'' es $f$.

\begin{definicion}[label=def:pm-ck]{Continuidad y clases $\Cont k$}
Sea $f:\R^n\to\R^m$, $f(x)=[f_1(x),\dots,f_m(x)]\T$.
\begin{itemize}
  \item $f$ es \clave{continua} en $x_0$ si $\lim_{x\to x_0}f(x)=f(x_0)$; equivalentemente, si cada $f_i$ es continua en $x_0$. Se escribe $f\in\Cont0$.
  \item $f\in\Cont1$ (\clave{continuamente diferenciable}) si todas las derivadas parciales $\partial f_i/\partial x_j$ existen y son continuas.
  \item $f\in\Cont k$ si las derivadas parciales hasta orden $k$ existen y son continuas; $f\in\Cont\infty$ (\clave{suave}) si las tiene de todos los órdenes.
\end{itemize}
\end{definicion}

Cada clase contiene a la siguiente: $\Cont\infty\subset\cdots\subset\Cont2\subset\Cont1\subset\Cont0$, y ninguna inclusión es una igualdad ($\abs x\in\Cont0\setminus\Cont1$; $x\abs x\in\Cont1\setminus\Cont2$).

\begin{cuidado}[Continua a trozos no quiere decir ``sin saltos'']
Una función $g(t)$ es \clave{continua a trozos} en $[t_0,t_1]$ si el intervalo se puede partir en un número finito de subintervalos en cuyo interior $g$ es continua y si en los extremos tiene límites laterales finitos. Puede tener \emph{saltos}: un escalón o una onda cuadrada son continuos a trozos. Lo que no puede tener son infinitos saltos ni explosiones. En el teorema de existencia se pide continuidad a trozos en $t$ precisamente para admitir entradas que conmutan.
\end{cuidado}

%% ---------------------------------------------------------------------
\subsection{El jacobiano}
%% ---------------------------------------------------------------------

\begin{definicion}[label=def:pm-jacobiano]{Jacobiano y gradiente}
El \clave{jacobiano} de $f:\R^n\to\R^m$ es la matriz $m\times n$ de sus derivadas parciales,
\[
  \jac fx=\begin{bmatrix}
    \partial f_1/\partial x_1 & \cdots & \partial f_1/\partial x_n\\
    \vdots & \ddots & \vdots\\
    \partial f_m/\partial x_1 & \cdots & \partial f_m/\partial x_n
  \end{bmatrix}:
  \quad\text{fila $i$ = componente $f_i$, columna $j$ = variable $x_j$.}
\]
Para una función escalar $h:\R^n\to\R$ el jacobiano es un vector \emph{fila}, $\jac hx\in\R^{1\times n}$, y el \clave{gradiente} es su transpuesta, $\nabla h(x)=\bigl(\jac hx\bigr)\T\in\R^n$.
\end{definicion}

El jacobiano es la mejor aproximación lineal de $f$ cerca de un punto: $f(x_0+\delta)=f(x_0)+\jacat fx{x_0}\delta+o(\norm\delta)$. Es la matriz que se calcula para linealizar en un equilibrio (capítulo~4). El gradiente apunta en la dirección de máximo crecimiento de $h$; la derivada de $h$ a lo largo de las soluciones de $\dot x=f(x)$ es $\dot h=\jac hx\,f(x)=\nabla h(x)\T f(x)$, la cuenta que repetiremos con cada función de Lyapunov.

\begin{ejemplo}[label=ej:pm-jacobiano]{Calcular un jacobiano}
Calcule el jacobiano de $f(x)=\begin{bmatrix}x_1^2+x_2x_3\\ x_1x_2+\sin x_3\\ \ee^{x_1}+x_2^2-x_3\end{bmatrix}$ y diga sus dimensiones.
\solucion Fila por componente, columna por variable:
\[
  \jac fx=\begin{bmatrix}2x_1&x_3&x_2\\ x_2&x_1&\cos x_3\\ \ee^{x_1}&2x_2&-1\end{bmatrix}\in\R^{3\times3} .
\]
Todas las entradas son continuas, así que $f\in\Cont1$ (de hecho $\Cont\infty$).
\end{ejemplo}

%% ---------------------------------------------------------------------
\subsection{El teorema del valor medio, en la forma correcta}
%% ---------------------------------------------------------------------

Para funciones escalares de una variable, $f(x)-f(y)=f'(z)(x-y)$ para algún $z$ entre $x$ y $y$. La versión que sirve para campos vectoriales es una integral.

\begin{lema}[label=lem:pm-valormedio]{Valor medio en forma integral}
Si $f\in\Cont1$ en un conjunto que contiene el segmento entre $x$ y $y$, entonces
\[
  f(x)-f(y)=\Bigl(\int_0^1\jacat fx{y+s(x-y)}\,\dif s\Bigr)(x-y) .
\]
\end{lema}

\begin{proof}
Sea $g(s)=f(y+s(x-y))$, $s\in[0,1]$. Por la regla de la cadena, $g'(s)=\jacat fx{y+s(x-y)}(x-y)$, y $f(x)-f(y)=g(1)-g(0)=\int_0^1g'(s)\,\dif s$.
\end{proof}

\begin{trampa}[label=tr:pm-valormedio]{No hay un punto intermedio para campos vectoriales}
\plausible{Como en una variable, $f(x)-f(y)=\jacat fx z(x-y)$ para algún $z$ en el segmento entre $x$ y $y$.}
Falso si $f$ tiene más de una componente. Tome $f(\theta)=[\cos\theta,\sin\theta]\T$ entre $\theta=0$ y $\theta=2\pi$: $f(2\pi)-f(0)=0$, pero $f'(z)\cdot2\pi=2\pi[-\sin z,\cos z]\T\neq0$ para todo $z$. Cada componente tiene su propio punto intermedio, y en general no coinciden.
\correcto{Use la forma integral (\cref{lem:pm-valormedio}) o, para acotar, la desigualdad $\norm{f(x)-f(y)}\le\sup_{z}\norm{\jacat fxz}\,\norm{x-y}$ con $z$ en el segmento. Esa desigualdad es la que convierte un jacobiano acotado en una constante de Lipschitz.}
\end{trampa}

%% ---------------------------------------------------------------------
\section{La condición de Lipschitz}
\label{sec:pm-lipschitz}
%% ---------------------------------------------------------------------

La continuidad dice que $f(x)$ se acerca a $f(y)$ cuando $x$ se acerca a $y$, pero no \emph{qué tan rápido}. La condición de Lipschitz lo cuantifica, y es exactamente lo que hace falta para la unicidad.

\begin{definicion}[label=def:pm-lipschitz]{Funciones Lipschitz}
Sea $f:[a,b]\times D\to\R^n$, con $D\subseteq\R^n$.
\begin{itemize}
  \item $f$ es \clave{Lipschitz en $x$ en un conjunto $W\subseteq D$} si existe $L\ge0$ tal que
  \begin{equation}
    \norm{f(t,x)-f(t,y)}\le L\norm{x-y}\qquad\text{para todo }x,y\in W,\ t\in[a,b] .
    \label{eq:pm-lipschitz}
  \end{equation}
  $L$ es una \clave{constante de Lipschitz} (cualquier número mayor también lo es).
  \item $f$ es \clave{localmente Lipschitz} en $D$ si cada punto de $D$ tiene una vecindad $W$ en la que $f$ es Lipschitz, con una constante que puede depender de la vecindad.
  \item $f$ es \clave{globalmente Lipschitz} si es Lipschitz en todo $\R^n$, con una sola constante.
\end{itemize}
\end{definicion}

Para $n=1$ e invariante en el tiempo, \eqref{eq:pm-lipschitz} dice $\abs{f(x)-f(y)}/\abs{x-y}\le L$: \emph{ninguna secante de la gráfica tiene pendiente mayor que $L$ en valor absoluto}. Geométricamente, en cada punto de la gráfica cabe un ``cono'' de pendiente $\pm L$ que contiene al resto de la gráfica (\cref{fig:pm-lipschitz}). El caso más simple es $f(x)=ax$: $\norm{ax-ay}=\abs a\norm{x-y}$, así que $L=\abs a$ sirve en todo $\R^n$.

\script{cap03_preliminares/C3_Ej5_Lipschitz.m}
\begin{figure}[ht]
\centering
\includegraphics[width=\linewidth]{cap03/f_lipschitz}
\caption{Conos de Lipschitz. $x^3$: en $[-a,a]$ cabe un cono de pendiente $3a^2$, pero no uno que sirva en toda la recta. $\sat$: un solo cono de pendiente $1$ sirve en todas partes, aunque $\sat$ no es derivable en $\pm1$. $x^{1/3}$: en el origen ningún cono sirve.}
\label{fig:pm-lipschitz}
\end{figure}

\begin{ejemplo}[label=ej:pm-tres]{Tres funciones, tres respuestas}
Clasifique $x^3$, $\sat(x)$ y $x^{1/3}$ en $\R$.
\solucion \emph{$x^3$: localmente, no globalmente.} $x^3-y^3=(x-y)(x^2+xy+y^2)$ y en $[-a,a]$ el segundo factor es a lo sumo $3a^2$; esa es la constante en $[-a,a]$, y crece sin cota con $a$. \emph{$\sat$: globalmente.} $\abs{\sat x-\sat y}\le\abs{x-y}$ para todo $x,y$ (por casos), con $L=1$, aunque $\sat\notin\Cont1$. \emph{$x^{1/3}$: ni siquiera localmente en $0$.} La secante entre $0$ y $h>0$ tiene pendiente $h^{1/3}/h=h^{-2/3}$, que crece sin cota cuando $h\to0$ (vale $100$ con $h=10^{-3}$ y $10^4$ con $h=10^{-6}$). Ningún $L$ sirve en ninguna vecindad del origen; fuera del origen sí es localmente Lipschitz.
\end{ejemplo}

\codigorepo{cap03_preliminares/C3_Ej5_Lipschitz.m}{secantes}{Pendientes de las secantes de $x^3$, $\sat$ y $x^{1/3}$}

%% ---------------------------------------------------------------------
\subsection{Del jacobiano a la constante de Lipschitz}
%% ---------------------------------------------------------------------

Calcular el supremo de los cocientes es poco práctico. Para funciones $\Cont1$ hay un atajo: acotar el jacobiano.

\begin{lema}[label=lem:pm-jacobiano]{Jacobiano acotado en un convexo \citep[Lema~3.1]{Kha02}}
Sea $f:[a,b]\times D\to\R^m$ continua, con $\partial f/\partial x$ continua en $[a,b]\times D$. Si $W\subseteq D$ es \emph{convexo} y $\norm{\partial f/\partial x(t,x)}\le L$ en $[a,b]\times W$ (norma inducida), entonces
\[
  \norm{f(t,x)-f(t,y)}\le L\norm{x-y}\qquad\text{para todo }t\in[a,b],\ x,y\in W .
\]
\end{lema}

\begin{proof}
El segmento entre $x$ y $y$ está en $W$ porque $W$ es convexo. Por el \cref{lem:pm-valormedio} y la \cref{prop:pm-matrices},
$\norm{f(t,x)-f(t,y)}\le\int_0^1\norm{\jacat fx{(t,y+s(x-y))}}\,\norm{x-y}\,\dif s\le L\norm{x-y}$.
\end{proof}

\begin{lema}[label=lem:pm-c1lip]{$\Cont1$ implica localmente Lipschitz \citep[Lema~3.2]{Kha02}}
Si $f$ y $\partial f/\partial x$ son continuas en $[a,b]\times D$, con $D$ un dominio, entonces $f$ es localmente Lipschitz en $x$ en $[a,b]\times D$.
\end{lema}

\begin{proof}
Cada $x_0\in D$ tiene una bola cerrada $W=\{\norm{x-x_0}\le r\}\subset D$. $W$ es convexa y compacta, y $\norm{\partial f/\partial x}$ es continua en el compacto $[a,b]\times W$: alcanza un máximo $L$. El \cref{lem:pm-jacobiano} hace el resto.
\end{proof}

\begin{lema}[label=lem:pm-global]{Lipschitz global \citep[Lema~3.3]{Kha02}}
Si $f$ y $\partial f/\partial x$ son continuas en $[a,b]\times\R^n$, entonces $f$ es globalmente Lipschitz en $x$ si y solo si $\partial f/\partial x$ está uniformemente acotada en $[a,b]\times\R^n$.
\end{lema}

\begin{proof}
Si $\norm{\partial f/\partial x}\le L$ en todas partes, el \cref{lem:pm-jacobiano} con $W=\R^n$ da $L$. Al revés, si $f$ es globalmente Lipschitz con constante $L$, cada derivada direccional es un límite de cocientes acotados por $L$: $\norm{\jacat fx{(t,x)}v}=\lim_{h\to0}\norm{f(t,x+hv)-f(t,x)}/\abs h\le L\norm v$, luego $\norm{\partial f/\partial x}\le L$.
\end{proof}

\begin{ejemplo}[label=ej:pm-lipcaja]{Una constante de Lipschitz en una caja}
Para $f(x)=\begin{bmatrix}-x_1+x_1x_2\\ x_2-x_1x_2\end{bmatrix}$ en $\R^2$: (a) ¿es localmente Lipschitz?, ¿globalmente? (b) Halle una constante de Lipschitz en $W=\{\abs{x_1}\le a_1,\ \abs{x_2}\le a_2\}$ con la norma $\infty$.
\solucion (a) $\jac fx=\begin{bmatrix}-1+x_2&x_1\\-x_2&1-x_1\end{bmatrix}$ es continuo, así que $f\in\Cont1$ y es localmente Lipschitz (\cref{lem:pm-c1lip}). No es globalmente Lipschitz: las entradas crecen sin cota y el \cref{lem:pm-global} lo descarta. (b) $W$ es convexo. La norma $\infty$ inducida es la máxima suma por filas: $\abs{-1+x_2}+\abs{x_1}\le1+a_1+a_2$ y $\abs{x_2}+\abs{1-x_1}\le1+a_1+a_2$. Luego $L=1+a_1+a_2$. Con $a_1=1$, $a_2=2$, $L=4$, y el mayor cociente encontrado con $50\,000$ pares aleatorios es $3.57$.
\end{ejemplo}

\codigorepo{cap03_preliminares/C3_Ej5_Lipschitz.m}{caja}{Verificación de $L=1+a_1+a_2$ con pares aleatorios}

\begin{trampa}[label=tr:pm-localglobal]{Local en cada punto no es Lipschitz en el conjunto}
\plausible{$f$ es localmente Lipschitz en cada punto de $D$; por lo tanto es Lipschitz en $D$.}
La constante puede crecer al moverse por $D$: $x^3$ es localmente Lipschitz en todo $\R$ y no es Lipschitz en $\R$. Tampoco basta con que $D$ sea acotado: $\sqrt x$ es localmente Lipschitz en $(0,1)$ y no es Lipschitz en $(0,1)$, porque la constante explota cerca del borde.
\correcto{Localmente Lipschitz en $D$ implica Lipschitz en cada \emph{compacto} $W\subset D$ (con una constante que depende de $W$). Esa es la versión que usan los teoremas: en un compacto se puede elegir una sola $L$.}
\end{trampa}

\begin{definicion}[label=def:pm-unifcont]{Continuidad uniforme}
$f$ es \clave{uniformemente continua} en $W$ si para todo $\epsilon>0$ existe $\delta>0$ ---el mismo para todos los puntos--- tal que $\norm{x-y}<\delta$ implica $\norm{f(x)-f(y)}<\epsilon$.
\end{definicion}

Si $f$ es Lipschitz en $W$ con constante $L>0$, es uniformemente continua ($\delta=\epsilon/L$). La \cref{fig:pm-regularidad} resume las relaciones en $\R$: $\Cont1\Rightarrow$ localmente Lipschitz; Lipschitz $\Rightarrow$ localmente Lipschitz \emph{y} uniformemente continua; ninguna de las implicaciones se invierte. $x^{1/3}$ es uniformemente continua en $\R$ sin ser localmente Lipschitz en el origen; $x^3$ es $\Cont1$ sin ser uniformemente continua.

\begin{figure}[ht]
\centering
\includegraphics[width=0.92\linewidth]{cap03/f_regularidad}
\caption{Clases de regularidad para funciones $\R\to\R$, con un representante en cada región.}
\label{fig:pm-regularidad}
\end{figure}

\begin{profundiza}[Lipschitz fuera de las ecuaciones diferenciales]
La constante de Lipschitz de una red neuronal mide cuánto puede cambiar la salida cuando cambia la entrada; acotarla controla la sensibilidad y el sobreajuste. Para una red residual, cada bloque es $x\mapsto x+\sigma(Wx+b)$, y con una activación $1$-Lipschitz su constante es a lo sumo $1+\norm W_2$; la de la red es a lo sumo el producto de las constantes de los bloques (composición y submultiplicatividad, \cref{prop:pm-matrices}). Con exactamente estas tres herramientas ---Lipschitz, normas matriciales y composición---, \citet{CGR26} diseñaron una regularización adaptativa para redes ResNet. Las redes residuales son, además, la discretización de Euler de una ecuación diferencial $\dot x=\sigma(W(t)x+b(t))$, y la condición de Lipschitz es la misma del teorema de existencia y unicidad.
\end{profundiza}

%% ---------------------------------------------------------------------
\section{Existencia y unicidad de soluciones}
\label{sec:pm-existencia}
%% ---------------------------------------------------------------------

Ya tenemos el lenguaje para responder la primera pregunta del capítulo. Una \emph{solución} de \eqref{eq:pm-pvi} en $[t_0,t_1]$ es una función continua $x:[t_0,t_1]\to\R^n$ tal que $\dot x(t)$ existe y vale $f(t,x(t))$ salvo, a lo sumo, en los instantes donde $f$ salta en $t$. Integrando, $x$ es solución si y solo si
\begin{equation}
  x(t)=x_0+\int_{t_0}^tf\bigl(s,x(s)\bigr)\,\dif s,\qquad t\in[t_0,t_1] .
  \label{eq:pm-integral}
\end{equation}
La ecuación integral tiene dos ventajas: no pide derivar (solo integrar, lo que tolera saltos en $t$) y convierte el problema en buscar un \emph{punto fijo}: si $(Px)(t)\defeq x_0+\int_{t_0}^tf(s,x(s))\,\dif s$, las soluciones son exactamente las funciones con $x=Px$.

%% ---------------------------------------------------------------------
\subsection{Dos herramientas}
%% ---------------------------------------------------------------------

\begin{lema}[label=lem:pm-gronwall]{Desigualdad de Gronwall--Bellman \citep[Lema~A.1]{Kha02}}
Sean $\lambda:[a,b]\to\R$ continua y $\mu:[a,b]\to\R$ continua y no negativa. Si una función continua $y:[a,b]\to\R$ cumple
\[
  y(t)\le\lambda(t)+\int_a^t\mu(s)\,y(s)\,\dif s\qquad\text{para }a\le t\le b,
\]
entonces, en el mismo intervalo,
\[
  y(t)\le\lambda(t)+\int_a^t\lambda(s)\,\mu(s)\exp\Bigl[\int_s^t\mu(\tau)\,\dif\tau\Bigr]\dif s .
\]
Si $\lambda\equiv\lambda_0$ es constante, $y(t)\le\lambda_0\exp\bigl[\int_a^t\mu(\tau)\,\dif\tau\bigr]$; si además $\mu\equiv\mu_0$, $y(t)\le\lambda_0e^{\mu_0(t-a)}$.
\end{lema}

\begin{proof}
Sea $z(t)=\int_a^t\mu(s)y(s)\,\dif s$, que es derivable con $\dot z=\mu y\le\mu(\lambda+z)$ (porque $\mu\ge0$ y $y\le\lambda+z$). Multiplicando por el factor integrante $\exp[-\int_a^t\mu]$,
\[
  \frac{\dif}{\dif t}\Bigl(z(t)e^{-\int_a^t\mu}\Bigr)=\bigl(\dot z-\mu z\bigr)e^{-\int_a^t\mu}\le\lambda\mu\,e^{-\int_a^t\mu}.
\]
Integrando de $a$ a $t$ con $z(a)=0$ y multiplicando por $e^{\int_a^t\mu}$: $z(t)\le\int_a^t\lambda(s)\mu(s)e^{\int_s^t\mu}\,\dif s$. Como $y\le\lambda+z$, queda la desigualdad. Si $\lambda\equiv\lambda_0$, la integral es $\lambda_0\int_a^t\mu(s)e^{\int_s^t\mu}\dif s=\lambda_0\bigl(e^{\int_a^t\mu}-1\bigr)$, de donde sale la segunda forma.
\end{proof}

\begin{ideaclave}
Gronwall--Bellman dice: una cantidad que crece \emph{a lo sumo} proporcionalmente a sí misma crece \emph{a lo sumo} exponencialmente. Convierte una desigualdad integral, que no se sabe resolver, en una cota explícita.
\end{ideaclave}

\begin{teorema}[label=thm:pm-contraccion]{Principio de contracción \citep{Ban22}, \citep[Teorema~B.1]{Kha02}}
Sea $S$ un subconjunto cerrado de un espacio de Banach $X$ y $T:S\to S$ un mapa con
\[
  \norm{T(x)-T(y)}\le\rho\norm{x-y}\qquad\text{para todo }x,y\in S,\qquad 0\le\rho<1 .
\]
Entonces existe un único $x^*\in S$ con $x^*=T(x^*)$, y las iteraciones $x_{k+1}=T(x_k)$ convergen a $x^*$ desde cualquier $x_1\in S$.
\end{teorema}

\begin{proof}
\emph{Cauchy.} $\norm{x_{k+1}-x_k}\le\rho\norm{x_k-x_{k-1}}\le\dots\le\rho^{k-1}\norm{x_2-x_1}$. Para $r\ge1$, por la triangular y la serie geométrica,
$\norm{x_{k+r}-x_k}\le\sum_{i=k}^{k+r-1}\rho^{i-1}\norm{x_2-x_1}\le\frac{\rho^{k-1}}{1-\rho}\norm{x_2-x_1}\to0$.
\emph{Límite en $S$.} $X$ es completo, así que $x_k\to x^*\in X$; $S$ es cerrado, así que $x^*\in S$.
\emph{Punto fijo.} $\norm{x^*-T(x^*)}\le\norm{x^*-x_{k+1}}+\norm{T(x_k)-T(x^*)}\le\norm{x^*-x_{k+1}}+\rho\norm{x_k-x^*}\to0$.
\emph{Unicidad.} Si $x^*$ y $y^*$ son puntos fijos, $\norm{x^*-y^*}=\norm{T(x^*)-T(y^*)}\le\rho\norm{x^*-y^*}$ con $\rho<1$; luego $x^*=y^*$.
\end{proof}

Las dos hipótesis que se usan de verdad son la completitud de $X$ (para que el límite exista) y que $S$ sea cerrado (para que el límite no se escape de $S$). Es aquí donde la \cref{sec:pm-completos} deja de ser abstracta.

%% ---------------------------------------------------------------------
\subsection{El teorema local}
%% ---------------------------------------------------------------------

\begin{teorema}[label=thm:pm-local]{Existencia y unicidad local \citep[Teorema~3.1]{Kha02}}
Sea $f(t,x)$ continua a trozos en $t$ y tal que
\[
  \norm{f(t,x)-f(t,y)}\le L\norm{x-y}\qquad\text{para todo }x,y\in B=\{x:\norm{x-x_0}\le r\},\ t\in[t_0,t_1] .
\]
Entonces existe $\delta>0$ tal que \eqref{eq:pm-pvi} tiene una solución única en $[t_0,t_0+\delta]$. Basta tomar
\begin{equation}
  \delta\le\min\Bigl\{t_1-t_0,\ \frac{r}{Lr+h},\ \frac\rho L\Bigr\},\qquad h=\max_{t\in[t_0,t_1]}\norm{f(t,x_0)},\quad \rho<1 .
  \label{eq:pm-delta}
\end{equation}
\end{teorema}

\begin{proof}
\emph{El espacio.} Sea $X=C[t_0,t_0+\delta]$ con $\norm\cdot_C$ (Banach, \cref{sec:pm-completos}) y $S=\{x\in X:\norm{x-x_0}_C\le r\}$, la bola de las funciones que no salen de $B$; $S$ es cerrado. Las soluciones de \eqref{eq:pm-pvi} son los puntos fijos de $P$.

\emph{$P$ manda $S$ en $S$.} Para $x\in S$, $f(s,x(s))-f(s,x_0)$ está acotado por $L\norm{x(s)-x_0}\le Lr$, y $\norm{f(s,x_0)}\le h$. Entonces
$\norm{(Px)(t)-x_0}\le\int_{t_0}^t\bigl(\norm{f(s,x(s))-f(s,x_0)}+\norm{f(s,x_0)}\bigr)\dif s\le(t-t_0)(Lr+h)\le\delta(Lr+h)\le r$.

\emph{$P$ es una contracción en $S$.} Para $x,y\in S$,
$\norm{(Px)(t)-(Py)(t)}\le\int_{t_0}^tL\norm{x(s)-y(s)}\,\dif s\le L\delta\norm{x-y}_C\le\rho\norm{x-y}_C$.
Por el \cref{thm:pm-contraccion}, $P$ tiene un único punto fijo en $S$: existe una única solución que se queda en $B$.

\emph{No hay otras soluciones fuera de $S$.} Sea $x$ cualquier solución en $[t_0,t_0+\delta]$. Si saliera de $B$, habría un primer instante $t_0+\mu$, $\mu\le\delta$, con $\norm{x(t_0+\mu)-x_0}=r$, y hasta ese instante estaría en $B$. Allí, con Gronwall--Bellman aplicado a $\norm{x(t)-x_0}\le h(t-t_0)+\int_{t_0}^tL\norm{x(s)-x_0}\,\dif s$, se obtiene $r=\norm{x(t_0+\mu)-x_0}\le\frac hL\bigl(e^{L\mu}-1\bigr)$, es decir $\mu\ge\frac1L\ln\bigl(1+\frac{Lr}h\bigr)>\frac{r}{Lr+h}\ge\delta$ (porque $\ln(1+u)>u/(1+u)$), una contradicción. Luego toda solución está en $S$ y es la del punto fijo.
\end{proof}

\notamargen{La palabra clave es \emph{local}: el teorema solo garantiza un intervalo $[t_0,t_0+\delta]$, y $\delta$ depende de $x_0$ a través de $h$ y de $r$.}
La prueba es constructiva: las iteraciones de Picard $x_{k+1}=Px_k$, desde $x_0(t)\equiv x_0$, convergen a la solución en la norma del máximo. Para $\dot x=x$, $x(0)=1$, en $[0,1]$, la iteración $k$ es el polinomio de Taylor $\sum_{j=0}^kt^j/j!$ de $e^t$, y el error en la norma del máximo cae como $e/(k+1)!$: $1.6\times10^{-3}$ en la quinta iteración (\cref{fig:pm-picard}).

\video{V3_1_picard}{Las iteraciones de Picard se pegan a $e^t$}
\begin{figure}[ht]
\centering
\includegraphics[width=\linewidth]{cap03/f_picard}
\caption{Aproximaciones sucesivas de Picard para $\dot x=x$, $x(0)=1$. Cada iteración aplica una vez el operador $P$; el error en la norma del máximo cae factorialmente.}
\label{fig:pm-picard}
\end{figure}

\codigorepo{cap03_preliminares/C3_Ej7_Picard.m}{picard}{Iteraciones de Picard: $x_{k+1}=x_0+\int_0^tf(x_k)$}

\begin{ejemplo}[label=ej:pm-delta]{¿Cuánto garantiza el teorema?}
Para $\dot x=x^2$, $x(0)=1$, calcule el mayor $\delta$ que da \eqref{eq:pm-delta} y compárelo con el intervalo donde realmente existe la solución.
\solucion En $B=\{\abs{x-1}\le r\}$, $\abs{x^2-y^2}=\abs{x+y}\abs{x-y}\le2(1+r)\abs{x-y}$: $L=2(1+r)$, y $h=\abs{f(1)}=1$. El cociente $r/(Lr+h)=r/(2r^2+2r+1)$ es máximo en $r^*=1/\sqrt2$, donde vale $(\sqrt2-1)/2\approx0.207$; el término $\rho/L$ con $\rho\to1$ vale allí $0.293$, mayor. El teorema garantiza $[0,\,0.207]$. La solución exacta, $x(t)=1/(1-t)$, existe en $[0,1)$ y escapa en $t=1$ (\cref{fig:pm-escape}, izquierda). El teorema es conservador, pero no es falso: \emph{nunca} promete más de lo que hay.
\end{ejemplo}

\codigorepo{cap03_preliminares/C3_Ej7_Picard.m}{delta}{El intervalo garantizado para $\dot x=x^2$}

%% ---------------------------------------------------------------------
\subsection{Cuando falla Lipschitz: no unicidad}
%% ---------------------------------------------------------------------

\begin{ejemplo}[label=ej:pm-nounicidad]{Infinitas soluciones desde el mismo punto}
Muestre que $\dot x=x^{1/3}$, $x(0)=0$, tiene infinitas soluciones y diga qué hipótesis del \cref{thm:pm-local} falla.
\solucion Para cada $c\ge0$, $x_c(t)=0$ si $t\le c$ y $x_c(t)=\bigl(\tfrac23(t-c)\bigr)^{3/2}$ si $t\ge c$ es solución: en $t>c$, $\dot x_c=\bigl(\tfrac23(t-c)\bigr)^{1/2}=x_c^{1/3}$, y las dos ramas empalman con derivada cero en $t=c$. La solución ``espera'' en el origen un tiempo arbitrario y luego despega (\cref{fig:pm-nounicidad}). Falla Lipschitz en el origen (\cref{ej:pm-tres}): la pendiente $\frac13x^{-2/3}$ es infinita en $0$. $f$ es continua, y por eso sí hay soluciones (teorema de Peano \citep{Pea1890}); lo que se pierde es la unicidad.
\end{ejemplo}

\video{V3_2_no_unicidad}{Todas parten de $x(0)=0$}
\begin{figure}[ht]
\centering
\includegraphics[width=0.78\linewidth]{cap03/f_no_unicidad}
\caption{Soluciones de $\dot x=x^{1/3}$ con $x(0)=0$. \cmd{ode45} devuelve $x\equiv0$ sin avisar que hay otras.}
\label{fig:pm-nounicidad}
\end{figure}

\codigorepo{cap03_preliminares/C3_Ej6_NoUnicidad.m}{ode45}{Lo que hace \cmd{ode45} desde $0$ y desde $10^{-12}$}

\begin{trampa}[label=tr:pm-simulador]{El simulador no detecta la no unicidad}
\plausible{Simulé $\dot x=x^{1/3}$ desde $x(0)=0$ con \cmd{ode45} y obtuve $x(t)\equiv0$; la solución es única y es cero.}
El integrador evalúa $f(0)=0$, da pasos de tamaño cero en $x$ y se queda en la rama $x\equiv0$. Desde $x(0)=10^{-12}$, que en cualquier experimento físico es indistinguible de cero, sigue la rama que despega y llega a $2\sqrt2\approx2.83$ en $t=3$. La simulación \emph{elige} una solución; no prueba que sea la única.
\correcto{La unicidad es una propiedad de $f$, y se verifica con la condición de Lipschitz, no con una simulación. Cuando falla, el modelo predice más de un futuro desde el mismo estado y hay que decidir, con información física adicional, cuál es el relevante.}
\end{trampa}

%% ---------------------------------------------------------------------
\subsection{¿Hasta cuándo existe la solución?}
%% ---------------------------------------------------------------------

Si $f$ es localmente Lipschitz en un dominio $D$, el \cref{thm:pm-local} se puede aplicar de nuevo en $t_0+\delta$, y otra vez, y así la solución se extiende a un \clave{intervalo máximo} $[t_0,T)$. Si $T<\infty$, la solución tiene que salir de todo compacto de $D$ cuando $t\to T$; en $D=\R^n$ eso quiere decir $\norm{x(t)}\to\infty$: \clave{escape en tiempo finito}. Para garantizar existencia en todo $t\ge t_0$ hay dos caminos.

\begin{teorema}[label=thm:pm-global]{Existencia global con Lipschitz global \citep[Teorema~3.2]{Kha02}}
Si $f(t,x)$ es continua a trozos en $t$ y existen $L$ y $h$ con
\[
  \norm{f(t,x)-f(t,y)}\le L\norm{x-y},\qquad\norm{f(t,x_0)}\le h\qquad\text{para todo }x,y\in\R^n,\ t\in[t_0,t_1],
\]
entonces \eqref{eq:pm-pvi} tiene una solución única en todo $[t_0,t_1]$.
\end{teorema}

\noindent La idea de la prueba: con Lipschitz global se puede tomar $r$ tan grande como se quiera en \eqref{eq:pm-delta}, y entonces $\delta$ no depende del punto de partida; se avanza en pasos de igual longitud hasta cubrir $[t_0,t_1]$.

\begin{teorema}[label=thm:pm-compacto]{Existencia global con un compacto \citep[Teorema~3.3]{Kha02}}
Sea $f(t,x)$ continua a trozos en $t$ y localmente Lipschitz en $x$ para $t\ge t_0$ y $x\in D$. Sea $W\subset D$ compacto con $x_0\in W$. Si se sabe que toda solución de \eqref{eq:pm-pvi} se queda en $W$ mientras existe, entonces existe una única solución definida para todo $t\ge t_0$.
\end{teorema}

\begin{proof}
En el compacto $W$, $f$ es Lipschitz con una sola constante (\cref{tr:pm-localglobal}). Si el intervalo máximo fuera $[t_0,T)$ con $T<\infty$, la solución tendría que salir de todo compacto de $D$ cuando $t\to T$, pero se queda en $W$.
\end{proof}

La gracia del \cref{thm:pm-compacto} es que se puede verificar \emph{sin resolver la ecuación}: basta mostrar que algún conjunto compacto es positivamente invariante. Para $\dot x=-x^3$, el intervalo $W=[-\abs{x_0},\abs{x_0}]$ es invariante porque $x\dot x=-x^4\le0$ (el tamaño $\abs x$ nunca crece). Las funciones de Lyapunov del \cref{cap:lyapunov} son, precisamente, una fábrica de compactos invariantes.

\script{cap03_preliminares/C3_Ej8_Escape.m}
\begin{figure}[ht]
\centering
\includegraphics[width=\linewidth]{cap03/f_escape}
\caption{Izquierda: $\dot x=x^2$ escapa en $t=1$; en verde, el intervalo que garantiza el \cref{thm:pm-local}. Derecha: $\dot x=-x^3$ tampoco es globalmente Lipschitz y su solución $x_0/\sqrt{1+2x_0^2t}$ existe para todo $t\ge0$.}
\label{fig:pm-escape}
\end{figure}

\begin{trampa}[label=tr:pm-global]{Lipschitz global es suficiente, no necesario}
\plausible{$f(x)=-x^3$ no es globalmente Lipschitz (su derivada $-3x^2$ no es acotada); por lo tanto sus soluciones pueden escapar en tiempo finito.}
El \cref{thm:pm-global} da una condición \emph{suficiente}. Que no se cumpla no dice nada. Aquí el término $-x^3$ empuja hacia el origen, y el \cref{thm:pm-compacto} con $W=[-\abs{x_0},\abs{x_0}]$ da existencia global; la solución explícita es $x(t)=x_0/\sqrt{1+2x_0^2t}$ (con $x_0=2$, $x(10)=2/9$). Lo que decide si hay escape no es el tamaño de $\abs{f'}$, sino hacia dónde apunta $f$: $\dot x=x^2$ escapa y $\dot x=-x^2$ (para $x_0>0$) no.
\correcto{Casi ningún sistema no lineal interesante es globalmente Lipschitz. La existencia global se prueba, en la práctica, encontrando un compacto positivamente invariante.}
\end{trampa}

\codigorepo{cap03_preliminares/C3_Ej8_Escape.m}{escape}{Escape de $\dot x=x^2$ y existencia global de $\dot x=-x^3$}

%% ---------------------------------------------------------------------
\subsection{Dependencia continua: ¿importan los errores de modelado?}
%% ---------------------------------------------------------------------

Un modelo nunca es exacto y una condición inicial nunca se conoce con precisión. Para que el modelo sirva, errores pequeños deben producir soluciones cercanas, al menos en intervalos finitos.

\begin{teorema}[label=thm:pm-dependencia]{Dependencia continua \citep[Teorema~3.4]{Kha02}}
Sea $f(t,x)$ continua a trozos en $t$ y Lipschitz en $x$ con constante $L$ en $[t_0,t_1]\times W$. Sean $y$ y $z$ soluciones de
\[
  \dot y=f(t,y),\ y(t_0)=y_0,\qquad\dot z=f(t,z)+g(t,z),\ z(t_0)=z_0,
\]
que se quedan en $W$ en $[t_0,t_1]$, con $\norm{g(t,x)}\le\mu$ en $[t_0,t_1]\times W$. Entonces, con $\gamma=\norm{y_0-z_0}$,
\begin{equation}
  \norm{y(t)-z(t)}\le\gamma\,e^{L(t-t_0)}+\frac\mu L\bigl(e^{L(t-t_0)}-1\bigr),\qquad t\in[t_0,t_1] .
  \label{eq:pm-dependencia}
\end{equation}
\end{teorema}

\begin{proof}
Restando las ecuaciones integrales y usando Lipschitz y la cota de $g$,
$\norm{y(t)-z(t)}\le\gamma+\mu(t-t_0)+\int_{t_0}^tL\norm{y(s)-z(s)}\,\dif s$.
Gronwall--Bellman con $\lambda(t)=\gamma+\mu(t-t_0)$ y $\mu(\cdot)\equiv L$ da
$\norm{y-z}\le\gamma+\mu(t-t_0)+\int_{t_0}^tL\bigl(\gamma+\mu(s-t_0)\bigr)e^{L(t-s)}\,\dif s$, y la integral (por partes) es $\gamma(e^{L(t-t_0)}-1)+\frac\mu L(e^{L(t-t_0)}-1)-\mu(t-t_0)$.
\end{proof}

\script{cap03_preliminares/C3_Ej10_Dependencia.m}
\begin{figure}[ht]
\centering
\includegraphics[width=0.78\linewidth]{cap03/f_dependencia}
\caption{$\dot y=-y+\tfrac12\sin y$ contra la versión perturbada con $\mu=0.1$ y $\gamma=0.05$ ($L=1.5$). La diferencia real se queda en $0.09$; la cota de Gronwall llega a $211$ en $t=5$.}
\label{fig:pm-dependencia}
\end{figure}

\begin{trampa}[label=tr:pm-cotaexp]{La cota exponencial no predice que las soluciones se separan}
\plausible{Por \eqref{eq:pm-dependencia}, dos soluciones que empiezan cerca se separan exponencialmente.}
\eqref{eq:pm-dependencia} es una \emph{cota superior}, válida para cualquier $f$ con constante $L$, incluso las que contraen. Para $\dot y=-y+\tfrac12\sin y$ la diferencia real con la versión perturbada nunca pasa de $0.091$ en $[0,5]$, mientras la cota llega a $211$ (\cref{fig:pm-dependencia}). El teorema solo usa $\norm{f(y)-f(z)}\le L\norm{y-z}$, que no distingue entre un sistema que acerca las soluciones y uno que las aleja.
\correcto{El teorema garantiza que en un intervalo \emph{finito} el error se puede hacer tan pequeño como se quiera achicando $\gamma$ y $\mu$. Para cotas que no crezcan con $t$ hace falta estabilidad, que es el tema de los capítulos~7 a~10.}
\end{trampa}

\codigorepo{cap03_preliminares/C3_Ej10_Dependencia.m}{dependencia}{La cota de Gronwall contra la diferencia real}

\begin{ideaclave}
\begin{tabular}{@{}ll@{}}
$f$ continua & existe una solución, quizá no única (Peano)\\
$f$ localmente Lipschitz & solución única en un intervalo máximo $[t_0,T)$\\
\quad\dots y una solución acotada & $T=\infty$ (\cref{thm:pm-compacto})\\
$f$ globalmente Lipschitz & solución única para todo $t$ (\cref{thm:pm-global})\\
cualquiera de los anteriores con Lipschitz & dependencia continua en intervalos finitos
\end{tabular}
\end{ideaclave}

\begin{historia}[Cauchy, Lipschitz, Peano, Picard, Lindelöf]
Cauchy probó existencia y unicidad en la década de 1820 para $f$ con derivada continua, con un método de poligonales. Lipschitz, en 1876, reemplazó la derivada por la condición que hoy lleva su nombre \citep{Lip1876}. Peano mostró en 1890 que la continuidad basta para la existencia \citep{Pea1890}. Picard introdujo ese mismo año las aproximaciones sucesivas \citep{Pic1890}, y Lindelöf, en 1894, dio la versión con la que se enuncia hoy el teorema \citep{Lin1894}. Por eso en muchos textos el \cref{thm:pm-local} se llama teorema de Picard--Lindelöf.
\end{historia}

%% ---------------------------------------------------------------------
\section{El lema de comparación}
\label{sec:pm-comparacion}
%% ---------------------------------------------------------------------

Sabemos cuándo existe una solución única, pero casi nunca la podemos escribir. La tercera pregunta del capítulo es si al menos se puede \emph{acotar}. La idea: si una cantidad escalar $v(t)$ ---una norma, una energía--- crece a lo sumo como la solución de una ecuación escalar sencilla, entonces se queda por debajo de esa solución.

Para enunciarlo con generalidad hace falta una derivada que exista aunque $v$ no sea derivable, porque $v(t)=\norm{x(t)}$ no lo es cuando $x(t)=0$. La \clave{derivada superior por la derecha} (de Dini) de una función continua $v$ es
\[
  D^+v(t)=\limsup_{h\to0^+}\frac{v(t+h)-v(t)}{h} .
\]
Si $v$ es derivable en $t$, $D^+v(t)=\dot v(t)$.

\begin{lema}[label=lem:pm-comparacion]{Lema de comparación \citep[Lema~3.4]{Kha02}}
Considere $\dot u=f(t,u)$, $u(t_0)=u_0$, con $u\in\R$, $f$ continua en $t$ y localmente Lipschitz en $u$ para $t\ge t_0$ y $u\in J\subseteq\R$. Sea $[t_0,T)$ el intervalo máximo de existencia de $u(t)$ ($T$ puede ser $\infty$), con $u(t)\in J$. Si $v(t)$ es continua, $v(t)\in J$, y
\[
  D^+v(t)\le f\bigl(t,v(t)\bigr),\qquad v(t_0)\le u_0,
\]
entonces $v(t)\le u(t)$ para todo $t\in[t_0,T)$.
\end{lema}

\noindent Dicho con palabras: si $v$ empieza por debajo de $u$ y en cada instante crece a lo sumo como crecería $u$ \emph{en el mismo valor}, no la puede adelantar. La hipótesis de Lipschitz no es decorativa: garantiza que $u$ es la única solución, y por lo tanto que no hay otra solución ``más alta'' con la que $v$ pudiera compararse. La prueba \citep[\S3.4]{Kha02} compara $v$ con las soluciones de $\dot u=f(t,u)+\lambda$, $\lambda>0$, y hace tender $\lambda\to0$ usando la dependencia continua.

\begin{ejemplo}[label=ej:pm-comparacion1]{Una cota sin resolver la ecuación}
Acote la solución de $\dot x=-(1+x^2)x$, $x(0)=a$, sin resolverla.
\solucion $f$ es $\Cont1$, así que hay una solución única en algún $[0,t_1)$. Sea $v=x^2$, que es derivable:
\[
  \dot v=2x\dot x=-2x^2-2x^4\le-2v,\qquad v(0)=a^2 .
\]
La ecuación de comparación $\dot u=-2u$, $u(0)=a^2$, tiene solución $u(t)=a^2e^{-2t}$, definida para todo $t$. Por el lema, $x(t)^2\le a^2e^{-2t}$, es decir,
\[
  \abs{x(t)}\le\abs a\,e^{-t}\qquad\text{mientras la solución exista.}
\]
Como la solución queda acotada, el \cref{thm:pm-compacto} dice que existe para todo $t\ge0$: la cota da, de paso, la existencia global. Con $a=2$: $x(1)=0.348\le2e^{-1}=0.736$.
\end{ejemplo}

\codigorepo{cap03_preliminares/C3_Ej9_Comparacion.m}{escalar}{La cota $\abs a e^{-t}$ contra la solución}

\begin{ejemplo}[label=ej:pm-comparacion2]{Comparación con la norma: por qué la derivada de Dini}
Para
\[
  \dot x_1=-x_1+\frac{2x_2}{1+x_2^2},\qquad\dot x_2=-x_2+\frac{2x_1}{1+x_1^2},
\]
pruebe que $\norm{x(t)}_2\le e^{-t}\norm{x(0)}_2+\sqrt2\,(1-e^{-t})$.
\solucion Escriba $\dot x=-x+g(x)$ con $g=\bigl[\frac{2x_2}{1+x_2^2},\frac{2x_1}{1+x_1^2}\bigr]\T$. Como $\abs{2s/(1+s^2)}\le1$ para todo $s$, $\norm{g(x)}_2\le\sqrt2$. Sea $v=\norm x_2$. Donde $x\neq0$, $v$ es derivable y, por Cauchy--Schwarz,
\[
  \dot v=\frac{x\T\dot x}{\norm x}=\frac{-\norm x^2+x\T g(x)}{\norm x}\le-v+\norm{g(x)}\le-v+\sqrt2 .
\]
Donde $x=0$, $v$ no es derivable, pero $D^+v=\lim_{h\to0^+}\norm{x(t+h)}/h=\norm{\dot x(t)}=\norm{g(0)}=0\le-0+\sqrt2$. Así que $D^+v\le-v+\sqrt2$ en todas partes; la ecuación de comparación $\dot u=-u+\sqrt2$ da la cota. Note lo que la cota \emph{no} dice: no dice que $x\to0$. El sistema tiene un equilibrio en $(1,1)$, con norma exactamente $\sqrt2$, y desde $(3,-2)$ la solución va hacia él (\cref{fig:pm-comparacion}, derecha).
\end{ejemplo}

\script{cap03_preliminares/C3_Ej9_Comparacion.m}
\begin{figure}[ht]
\centering
\includegraphics[width=\linewidth]{cap03/f_comparacion}
\caption{Izquierda: la cota del \cref{ej:pm-comparacion1}. Derecha: la cota del \cref{ej:pm-comparacion2}; la solución se acerca al equilibrio $(1,1)$, cuya norma es el valor final $\sqrt2$ de la cota.}
\label{fig:pm-comparacion}
\end{figure}

\codigorepo{cap03_preliminares/C3_Ej9_Comparacion.m}{vectorial}{La cota de comparación para la norma}

\begin{trampa}[label=tr:pm-comparacion]{Comparar con la función equivocada}
\plausible{$\dot v\le-v+\sqrt2$ y $v(0)$ es pequeño; por lo tanto $v(t)\to0$.}
La comparación da $v\le u$, donde $u$ es la solución de la ecuación con \emph{igualdad}. Aquí $u(t)\to\sqrt2$, no $0$, y la conclusión correcta es $\limsup v\le\sqrt2$. En el \cref{ej:pm-comparacion2} el límite real es $\sqrt2$: la cota es ajustada y $x\not\to0$. El error es leer la desigualdad como si la ecuación de comparación fuera $\dot u=-u$.
\correcto{Resuelva la ecuación de comparación completa, con el mismo término que aparece en la desigualdad, y lea la cota que da. Una cota que no tiende a cero no prueba estabilidad, pero sí prueba que la solución es acotada y, con el \cref{thm:pm-compacto}, que existe para todo $t$.}
\end{trampa}

\begin{pausa}
Si $v\ge0$ cumple $\dot v\le-v^2$ con $v(0)=v_0>0$, ¿qué cota da el lema de comparación? ¿Converge $v$ a cero exponencialmente?
\respuesta{La ecuación de comparación $\dot u=-u^2$, $u(0)=v_0$, tiene solución $u=v_0/(1+v_0t)$. Luego $v(t)\le v_0/(1+v_0t)$: $v\to0$, pero solo como $1/t$, no exponencialmente. Es, salvo constantes, lo que pasa con $\dot x=-x^3$: con $v=x^2$, $\dot v=-2v^2$.}
\end{pausa}

%% ---------------------------------------------------------------------
\section{Para llevar}
\label{sec:pm-cierre}
%% ---------------------------------------------------------------------

\begin{parallevar}
\begin{enumerate}
  \item \textbf{Espacio $=$ conjunto $+$ estructura.} Una métrica mide distancias; una norma, tamaños; un producto interno, ángulos. Cada una induce la anterior. La completitud se agrega aparte: normado y completo es Banach; con producto interno y completo, Hilbert.
  \item \textbf{En $\R^n$ todas las normas son equivalentes}, así que acotado, convergente o estable no depende de la norma (las constantes sí). En espacios de señales las normas \emph{no} son equivalentes: hay que decir en cuál.
  \item \textbf{Normas de matrices:} $\norm A_1$ por columnas, $\norm A_\infty$ por filas, $\norm A_2=\sqrt{\lmax(A\T A)}$, que en general es mayor que el radio espectral.
  \item \textbf{Regularidad:} $\Cont1\Rightarrow$ localmente Lipschitz; con jacobiano acotado en un convexo, Lipschitz en ese convexo; con jacobiano uniformemente acotado en $\R^n$, globalmente Lipschitz.
  \item \textbf{Existencia y unicidad:} localmente Lipschitz da una solución única en un intervalo máximo. Se prueba con el principio de contracción en $(C[t_0,t_0+\delta],\norm\cdot_C)$, y la unicidad, con Gronwall--Bellman.
  \item \textbf{Existencia global:} con Lipschitz global, o ---lo que se usa en la práctica--- mostrando que la solución se queda en un compacto.
  \item \textbf{Gronwall y comparación:} una cantidad que crece a lo sumo como $f(t,v)$ se queda debajo de la solución de $\dot u=f(t,u)$. Con $v=\norm x$ hace falta la derivada de Dini.
\end{enumerate}
\end{parallevar}

\begin{trampa}[title={\faBug\hspace{0.45em}Las trampas del capítulo, en una caja}]{}
\begin{itemize}
  \item Parecer una distancia no basta: falta la triangular (\cref{tr:pm-parece}); y la fórmula de $\norm\cdot_p$ no es norma si $p<1$ (\cref{tr:pm-pmenor}).
  \item La equivalencia de normas es un privilegio de $\R^n$ (\cref{tr:pm-dimension}); $\norm A_2$ no es el mayor $\abs\lambda$ (\cref{tr:pm-radio}).
  \item Cauchy--Schwarz no sale de $\abs{\cos\theta}\le1$ (\cref{tr:pm-circular}); producto interno no es Hilbert sin completitud (\cref{tr:pm-completo}).
  \item El valor medio no tiene un punto intermedio para campos vectoriales (\cref{tr:pm-valormedio}); local en cada punto no es Lipschitz en el conjunto (\cref{tr:pm-localglobal}).
  \item El simulador no detecta la no unicidad (\cref{tr:pm-simulador}); Lipschitz global es suficiente, no necesario (\cref{tr:pm-global}).
  \item La cota de Gronwall no predice separación (\cref{tr:pm-cotaexp}); compare con la ecuación completa (\cref{tr:pm-comparacion}).
\end{itemize}
\end{trampa}

\begin{profundiza}[Correcciones respecto a las versiones anteriores (PrelimMath, Handout6a y las diapositivas de 2026)]
(i) Un espacio de Banach no es ``un conjunto con norma'' ni ``el conjunto unido con todas las normas posibles'': es un espacio normado \emph{completo}; lo mismo para Hilbert con producto interno. (ii) $(C[0,\infty),\mathcal{L}_p)$ no es un espacio de Banach: $C[a,b]$ con la norma $\mathcal{L}_2$ no es completo (\cref{tr:pm-completo}), y en $[0,\infty)$ ni siquiera toda función continua tiene norma finita. (iii) La desigualdad de Cauchy--Schwarz es un teorema, no una definición, y su prueba no puede usar $\cos\theta$ (\cref{tr:pm-circular}); se corrigió la prueba de $\norm x_\infty\le\norm x_2$, que tenía una errata. (iv) La equivalencia de normas no requiere $\alpha,\beta\in\mathbb N$: vale para \emph{cualquier} par de normas en $\R^n$. (v) Continua a trozos admite saltos; la versión anterior decía lo contrario. (vi) $\sat$ es \emph{globalmente} Lipschitz, no solo localmente. (vii) El resumen de existencia decía que con Lipschitz global la solución es única ``en un intervalo finito $[0,\infty)$''; el enunciado correcto es el \cref{thm:pm-global}, y la existencia global se obtiene en la práctica con el \cref{thm:pm-compacto}. (viii) En Gronwall--Bellman, la forma con $\lambda$ constante es $\lambda e^{\int_a^t\mu}$, no $\lambda e^{\int_s^t\mu}$; se completó la prueba y se agregó la parte de la unicidad que faltaba en la demostración del teorema local. (ix) En el diagrama de regularidad, Lipschitz está contenido en localmente Lipschitz (\cref{fig:pm-regularidad}); el diagrama anterior lo dibujaba parcialmente afuera. (x) Se agregaron las figuras que decían ``AÑADIR FIGURA'' (dominio y convexidad, \cref{fig:pm-conjuntos}). (xi) La notación se unificó con \texttt{notacion.sty}: sin negrillas, $f(t,x)$, $\inner xy$ en lugar de $<x,y>$, $\mathcal{L}_p$ para las señales. Todos los números se verificaron con \texttt{herramientas/cap03\_numeros.py} (43 de 43) y \cmd{verificarCap03.m} (33 de 33).
\end{profundiza}

%% ---------------------------------------------------------------------
\section{Problemas}
\label{sec:pm-problemas}
%% ---------------------------------------------------------------------

Los problemas marcados con \etiquetaM{} requieren \matlab; los marcados con \etiquetaEstrella{} son de profundización. Cada uno tiene una trampa: búsquela antes de dar la respuesta.

\subsection*{Espacios, normas y productos internos}

\begin{problema}[una métrica acotada]\label{prob:pm-acotada}
(a) Pruebe que $d(x,y)=\dfrac{\abs{x-y}}{1+\abs{x-y}}$ es una métrica en $\R$. (Pista: $s\mapsto s/(1+s)$ es creciente y subaditiva.) (b) Con esta métrica, ninguna distancia pasa de $1$. ¿Es $\R$ ``acotado'' con $d$? ¿Es compacto? (c) ¿Viene $d$ de alguna norma?
\end{problema}

\begin{problema}[normas pesadas]\label{prob:pm-pesada}
Sea $P\in\R^{n\times n}$. (a) Pruebe que $\norm x_P=\sqrt{x\T Px}$ es una norma si $P$ es simétrica y definida positiva, y que viene del producto interno $\inner xy_P=x\T Py$. (b) Muestre que $\lmin(P)\norm x_2^2\le x\T Px\le\lmax(P)\norm x_2^2$. (c) Para $P=\begin{bmatrix}1&2\\0&1\end{bmatrix}$, ¿es $\sqrt{x\T Px}$ una norma? (d) ¿Y para $P=\begin{bmatrix}1&0\\0&0\end{bmatrix}$?
\end{problema}

\begin{problema}[señales]\label{prob:pm-senales}
Diga a cuáles de $\mathcal{L}_1$, $\Ldos$ y $\Linf$ en $[0,\infty)$ pertenece cada señal y calcule las normas finitas: (a) $e^{-t}\cos t$; (b) $te^{-t}$; (c) $1/\sqrt{1+t}$; (d) $\sin(t)/(1+t)$.
\end{problema}

\begin{problema}[crecimiento transitorio]\label{prob:pm-transitorio}
Sea $A=\begin{bmatrix}-1&10\\0&-2\end{bmatrix}$. (a) ¿Es $A$ Hurwitz? (b) Calcule $\norm A_2$ y $\rho(A)$. (c) \etiquetaM{} Grafique $\norm{e^{At}}_2$ para $t\in[0,6]$ y halle su máximo. (d) Explique con la \cref{tr:pm-radio} por qué una matriz con valores propios $-1$ y $-2$ puede amplificar más de $2.5$ veces una condición inicial.
\end{problema}

\subsection*{Regularidad}

\begin{problema}[clasificar]\label{prob:pm-clasificar}
Clasifique en $\R$ (continua, $\Cont1$, localmente Lipschitz, globalmente Lipschitz) y dé una constante de Lipschitz cuando exista: (a) $\abs x$; (b) $x\abs x$; (c) $\sqrt{\abs x}$; (d) $\tanh x$; (e) $x+\sin x$. ¿Cuál de ellas es globalmente Lipschitz sin ser $\Cont1$?
\end{problema}

\begin{problema}[el péndulo]\label{prob:pm-pendulo}
Para $\dot x_1=x_2$, $\dot x_2=-\sin x_1-bx_2$, $b>0$: (a) calcule el jacobiano y muestre que $f$ es globalmente Lipschitz. (b) Dé una constante de Lipschitz con la norma $\infty$ y otra con la norma $2$. (c) Concluya que las soluciones existen para todo $t$ sin resolver nada. (d) ¿Sigue valiendo si la fricción es $-bx_2^3$?
\end{problema}

\subsection*{Existencia, unicidad y cotas}

\begin{problema}[el tanque que se vacía]\label{prob:pm-torricelli}
El nivel de un tanque que se vacía por un orificio cumple $\dot h=-k\sqrt h$, $h\ge0$. (a) Resuelva desde $h(0)=h_0>0$ y halle el instante $T$ en que se vacía. (b) Muestre que, hacia atrás en el tiempo, la solución que pasa por $h=0$ no es única. (c) ¿Qué hipótesis del \cref{thm:pm-local} falla y dónde? (d) Interprete físicamente.
\end{problema}

\begin{problema}[escape en $\pi/2$]\label{prob:pm-tan}
Para $\dot x=1+x^2$, $x(0)=0$: (a) resuelva y halle el tiempo de escape. (b) Calcule el $\delta$ que garantiza \eqref{eq:pm-delta}, optimizando en $r$. (c) \etiquetaM{} Implemente las iteraciones de Picard en $[0,\delta]$ y en $[0,1.5]$. ¿Convergen en los dos intervalos?
\end{problema}

\begin{problema}[un compacto invariante]\label{prob:pm-compacto}
Para $\dot x=x-x^3$: (a) ¿es globalmente Lipschitz? (b) Muestre que para todo $x_0$ el intervalo $W=[-\max\{1,\abs{x_0}\},\max\{1,\abs{x_0}\}]$ es positivamente invariante. (c) Concluya existencia global con el \cref{thm:pm-compacto}.
\end{problema}

\begin{problema}[comparación con entrada]\label{prob:pm-entrada}
Para $\dot x=-x-x^3+u(t)$ con $\abs{u(t)}\le1$, use $v=\abs x$ y la derivada de Dini para probar que $\abs{x(t)}\le e^{-t}\abs{x_0}+1-e^{-t}$. ¿Por qué no se puede usar $\dot v$ directamente? ¿Qué cota da si usa también el término $-x^3$? \etiquetaEstrella{}
\end{problema}

\begin{problema}[unicidad con Gronwall]\label{prob:pm-unicidadGB}
\etiquetaEstrella{} Sin usar el principio de contracción: si $x$ y $y$ son dos soluciones de \eqref{eq:pm-pvi} que se quedan en $B$, donde $f$ es Lipschitz con constante $L$, pruebe con Gronwall--Bellman que $x\equiv y$.
\end{problema}

%% ---------------------------------------------------------------------
\section{Soluciones seleccionadas}
%% ---------------------------------------------------------------------

\paragraph{Problema~\ref{prob:pm-acotada}}
(a) Con $\varphi(s)=s/(1+s)$, creciente y con $\varphi(s+t)\le\varphi(s)+\varphi(t)$: $d(x,y)=\varphi(\abs{x-y})\le\varphi(\abs{x-z}+\abs{z-y})\le d(x,z)+d(z,y)$. (b) Todo $\R$ está a distancia menor que $1$ de $0$: es acotado con $d$, pero no compacto (la sucesión $x_k=k$ no tiene subsucesión convergente). ``Cerrado y acotado implica compacto'' es un teorema sobre $\R^n$ \emph{con una norma}. (c) No: una métrica que viene de una norma cumple $d(\alpha x,0)=\abs\alpha d(x,0)$, y esta no.

\paragraph{Problema~\ref{prob:pm-pesada}}
(c) No: $x\T Px$ solo ve la parte simétrica $\tfrac12(P+P\T)=\begin{bmatrix}1&1\\1&1\end{bmatrix}$, que es semidefinida, y $x=[1,-1]\T$ da $x\T Px=1-2+1=0$ con $x\neq0$. (d) No: $x=[0,1]\T$ tiene ``norma'' cero. La trampa es mirar $P$ y no su parte simétrica.

\paragraph{Problema~\ref{prob:pm-transitorio}}
(a) Sí, valores propios $-1$ y $-2$. (b) $\rho=2$ y $\norm A_2\approx10.2$. (c) $e^{At}=\begin{bmatrix}e^{-t}&10(e^{-t}-e^{-2t})\\0&e^{-2t}\end{bmatrix}$; la entrada fuera de la diagonal llega a $2.5$ en $t=\ln2$, y $\norm{e^{At}}_2$ llega a $2.56$ en $t\approx0.66$ antes de decaer. (d) $A$ no es normal; los valores propios dicen hacia dónde converge, no cuánto amplifica en el camino.

\paragraph{Problema~\ref{prob:pm-clasificar}}
(a) Globalmente Lipschitz con $L=1$, no $\Cont1$. (b) $\Cont1$ (derivada $2\abs x$), localmente pero no globalmente Lipschitz. (c) Continua, no localmente Lipschitz en $0$ (como $x^{1/3}$). (d) $\Cont\infty$ y globalmente Lipschitz, $L=1$. (e) $\Cont\infty$ y globalmente Lipschitz, $L=2$. La respuesta a la pregunta final es (a).

\paragraph{Problema~\ref{prob:pm-torricelli}}
(a) $\sqrt{h(t)}=\sqrt{h_0}-kt/2$, el tanque se vacía en $T=2\sqrt{h_0}/k$ y después $h\equiv0$. (b) Desde $h(T)=0$ hacia atrás, cualquier $h_0$ (y cualquier instante de vaciado anterior) es compatible. (c) $\sqrt h$ no es Lipschitz en $h=0$. (d) Mirando un tanque vacío no se puede saber cuándo se vació: la información se pierde. Hacia adelante sí hay unicidad (el tanque vacío se queda vacío), porque $f\le0$ empuja contra el borde $h=0$.

\paragraph{Problema~\ref{prob:pm-tan}}
(a) $x=\tan t$, escape en $t=\pi/2\approx1.571$. (b) En $\{\abs x\le r\}$: $L=2r$, $h=1$; $r/(2r^2+1)$ es máximo en $r=1/\sqrt2$ con valor $1/(2\sqrt2)\approx0.354$, y $\rho/L<1/\sqrt2$: $\delta\approx0.354$. (c) Numéricamente convergen en los dos intervalos (con $60$ iteraciones el error en $[0,1.5]$ es del orden de $10^{-5}$), aunque el teorema solo lo garantiza en $[0,\delta]$: la conclusión del teorema es una cota inferior del intervalo, no superior.

\paragraph{Problema~\ref{prob:pm-entrada}}
Donde $x\neq0$, $\dot v=\sgn(x)\dot x=-\abs x-\abs x^3+\sgn(x)u\le-v+1$. Donde $x=0$, $\abs x$ no es derivable, pero $D^+v=\abs{\dot x}=\abs u\le1=-0+1$. La comparación con $\dot w=-w+1$ da la cota. Usando $-v^3$ se obtiene una cota mejor, pero la ecuación de comparación $\dot w=-w-w^3+1$ no tiene solución explícita.

\respuestaspausas

